\chapter{深度、正则性、对偶性}

本章中，所有环均假设为交换环，代数均假设为结合、交换且含单位元的，代数同态均保持单位元。若 A 是局部环，则以 \(m_{A}\) 表示其极大理想，以 \(k_{A}\) 表示其剩余域。若 p 是环 A 的素理想，则以 \(\kappa(\mathfrak{p})\) 表示局部环 \(A_{p}\) 的剩余域；把它与整环 A/p 的分式域恒同。以 \(\overline{Z}\) 表示子集 \(Z \cup \{-\infty, +\infty\}\) ，它含于 \(\overline{R}\) 中。于是，对一切 \(n \in Z\) ，有关系式 \(-\infty < n < +\infty\) 及 \(n + \infty = \infty + n = \infty + \infty = \infty\) ，\(n - \infty = -\infty + n = -\infty - \infty = -\infty\) 。

\section{深度}

\subsection{深度的同调定义}

设 A 为一个环，J 为 A 的一个理想，M 为一个 A-模。

定义 1. M 关于 J 的深度称为并记作 \(\mathrm{prof}_{\mathrm{A}}(\mathrm{J};\mathrm{M})\) ，或 \(\mathrm{prof}(\mathrm{J};\mathrm{M})\) ，即 \(\mathbf{N}\cup\{+\infty\}\) 中由满足下述条件的整数 \(n\) 组成的集合的下确界：\(\mathrm{Ext}_{\mathrm{A}}^{n}(\mathrm{A}/\mathrm{J},\mathrm{M})\) 非零。

当环 A 为局部环时，把 M 的深度简单地称为并记作 \(\text{prof}_{\mathrm{A}}(\mathrm{M})\) 或 \(\text{prof}(\mathrm{M})\) 的，就是 M 关于 A 的极大理想 \(m_{A}\) 的深度；而局部环 A 的深度则指 A-模 A 的深度。

\(\operatorname{Si} \operatorname{prof}_{\mathrm{A}}(\mathrm{J}; \mathrm{M}) = +\infty\) ，则 \(\operatorname{Ext}_{\mathrm{A}}^{i}(\mathrm{A}/\mathrm{J}, \mathrm{M}) = 0\) 对一切 \(i\) 成立。\(\operatorname{Si} \operatorname{prof}_{\mathrm{A}}(\mathrm{J}; \mathrm{M}) = r < +\infty\) ，则 \(\operatorname{Ext}_{\mathrm{A}}^{i}(\mathrm{A}/\mathrm{J}, \mathrm{M}) = 0\) 对 \(i < r\) 成立，且 \(\operatorname{Ext}_{\mathrm{A}}^{r}(\mathrm{A}/\mathrm{J}, \mathrm{M}) \neq 0\) 。

注.--- 1) 设 A-模 M 有限生成且 M = JM，即 Supp(M) ∩ V(J) = ∅ (II, § 4, 第 4 节, 命题 18 的推论)。此时 prof\_A(J;M) 等于 +∞：事实上，理想 Ann(M) + J 这时等于 A（同上引文），且对一切 i 含于 Ext\_A\^{}i(A/J;M) 的零化子之中。下文（第 3 节, 定理 1 的推论 1）将看到，反之，若理想 J 有限生成，则 prof\_A(J;M) = +∞ 蕴含 M = JM。

\begin{enumerate}
\def\labelenumi{\arabic{enumi})}
\setcounter{enumi}{1}
\tightlist
\item
  为使 \(\mathrm{prof}_{\mathbf{A}}(\mathbf{J};\mathbf{M})\) 为零，必须且只需 \(\mathrm{Hom}_{\mathbf{A}}(\mathbf{A} / \mathbf{J},\mathbf{M})\) 非零，即 M 具有一个被 J 零化的非零元素；特别地，当 \(\mathrm{Ass(M)}\cap \mathrm{V(J)}\neq \emptyset\) 时即属此情形。若环 A 为 Noether 环且 A-模 M 有限生成，则下列条件等价 (IV, § 1, 第 4 节, 命题 8)：
\end{enumerate}

\begin{enumerate}
\def\labelenumi{\alph{enumi})}
\item
  \(\mathrm{prof}_{\mathrm{A}}(\mathrm{J};\mathrm{M}) = 0\)
\item
  对每个 \(x \in J\) ，乘以 x 的映射 \(x_{\mathbb{M}}\) 都不是单射；
\end{enumerate}

\[
\mathrm{c)} \text {   we have   } \operatorname{Ass} (\mathrm{M}) \cap \mathrm{V} (\mathrm{J}) \neq \varnothing .
\]

\begin{enumerate}
\def\labelenumi{\arabic{enumi})}
\setcounter{enumi}{2}
\item
  根据注 2，为使局部环 \(\Lambda\) 的深度为零，必须且只需存在 A 的非零元素 \(x\) 使 \(\mathfrak{m}_{\mathrm{A}}x = 0\) 。若 A 不是域，则元素 \(x\) 一旦满足 \(\mathfrak{m}_{\mathrm{A}}x = 0\) 即不可逆，故属于 \(\mathfrak{m}_{\mathrm{A}}\) ，从而满足 \(x^{2} = 0\) 。因此，维数 \(\geqslant 1\) 的约化局部环的深度 \(\geqslant 1\) 。
\item
  设 \((\mathbf{M}_{\iota})_{\iota \in I}\) 为一族 A-模。根据 A, X, 第 89 页, 命题 7，有 \(\operatorname{prof}(\mathbf{J};\prod_{\iota \in I}\mathbf{M}_{\iota}) = \inf_{\iota \in I}\operatorname{prof}(\mathbf{J};\mathbf{M}_{\iota})\) 。
\end{enumerate}

命题 1.--- 设 A 为环，J 为 A 的理想，且 0 → M' → M → M'' → 0 为 A-模的正合列。令

\[
p ^ {\prime} = \operatorname{prof} (\mathrm{J}; \mathrm{M} ^ {\prime}), p = \operatorname{prof} (\mathrm{J}; \mathrm{M}), p ^ {\prime \prime} = \operatorname{prof} (\mathrm{J}; \mathrm{M} ^ {\prime \prime}).
\]

这时必处于以下三种两两互斥的情形之一：

\[
p ^ {\prime} = p \leqslant p ^ {\prime \prime}, \quad p = p ^ {\prime \prime} <   p ^ {\prime}, \quad p ^ {\prime \prime} = p ^ {\prime} - 1 <   p.
\]

考察与 A/J 及上述正合列相伴的扩张模正合列 (A, X, 第 92 页, 定理 2)。排除 \(p = p' = p'' = +\infty\) 的情形；此时该序列中存在第一个非零模，且紧随其后的模也非零。由此得到以下三种可能：

\begin{enumerate}
\def\labelenumi{\alph{enumi})}
\item
  第一个非零模是 \(\operatorname{Ext}_{\Lambda}^{p'}(A/J, M')\) 。于是 \(p' = p \leqslant p''\) 。
\item
  第一个非零模是 \(\operatorname{Ext}_{\Lambda}^{p}(\mathrm{A} / \mathrm{J},\mathrm{M})\) 。于是 \(p = p'' < p'\) 。
\item
  第一个非零模是 \(\mathrm{Ext}_{\mathrm{A}}^{p''}(\mathrm{A} / \mathrm{J},\mathrm{M}'')\) 。于是 \(p'' + 1 = p' \leqslant p\) 。
\end{enumerate}

注 5.--- 设 \(p = p'\) ，且出现在命题 1 的正合列中的单射 \(u: M' \to M\) 属于 J Hom \(_A(M', M)\) 。于是有 \(a\) \(p'' = p - 1\) 。事实上，该假设蕴含映射 Ext \(_A^i(1_{A/J}, u)\) 对每个整数 \(i\) 均为零；这就排除了上面所考虑的情形 \(a\) )。

命题 2.--- 设 A 为环，J 为 A 的理想，M 为 A-模，N 为被 J 的某个幂所零化的 A-模。则 \(\mathrm{Ext}_{\mathrm{A}}^{i}(\mathrm{N},\mathrm{M}) = 0\) 对每个整数 \(i < \mathrm{prof}_{\mathrm{A}}(\mathrm{J};\mathrm{M})\) 成立。

首先设 JN = 0，并对整数 \(i < \mathrm{prof}_{\mathrm{A}}(\mathrm{J};\mathrm{M})\) 作归纳论证。断言对 \(i < 0\) 显然。把 N 视为 (A/J)-模，并选取 (A/J)-模的一个正合列

\[
0 \rightarrow \mathrm{K} \longrightarrow (\mathrm{A} / \mathrm{J}) ^ {(\mathrm{I})} \longrightarrow \mathrm{N} \rightarrow 0.
\]

由此导出扩张模的一个正合列

\[
\operatorname{Ext} _ {\mathrm{A}} ^ {i - 1} (\mathrm{K}, \mathrm{M}) \longrightarrow \operatorname{Ext} _ {\mathrm{A}} ^ {i} (\mathrm{N}, \mathrm{M}) \longrightarrow \operatorname{Ext} _ {\mathrm{A}} ^ {i} ((\mathrm{A} / \mathrm{J}) ^ {(1)}, \mathrm{M}).
\]

A-模 \(\mathrm{Ext}_{\mathrm{A}}^{i-1}(\mathrm{K},\mathrm{M})\) 由归纳假设为零，而 A-模 \(\mathrm{Ext}_{\mathrm{A}}^{i}((\mathrm{A}/\mathrm{J})^{(1)},\mathrm{M})\) 同构于 \(\mathrm{Ext}_{\mathrm{A}}^{i}(\mathrm{A}/\mathrm{J},\mathrm{M})^{\mathrm{I}}\) (A, X, 第 89 页, 命题 7)，后者依深度的定义而为零。于是有 \(\mathrm{Ext}_{\mathrm{A}}^{i}(\mathrm{N},\mathrm{M}) = 0\) 。

转入一般情形，并对最小的整数 m \textgreater{} 0 作归纳论证，其中 m 满足 \(J^{m}N = 0\) 。我们刚才已处理 m = 1 的情形。设 m \textgreater{} 1 ，并设整数 \(i < \text{prof}_{\text{A}}(\text{J}; \text{M})\) 。考察正合列

\[
\operatorname{Ext} _ {\Lambda} ^ {i} (\mathrm{N/JN,M}) \longrightarrow \operatorname{Ext} _ {\Lambda} ^ {i} (\mathrm{N,M}) \longrightarrow \operatorname{Ext} _ {\Lambda} ^ {i} (\mathrm{JN,M})
\]

它由正合列 \(0 \rightarrow JN \rightarrow N \rightarrow N/JN \rightarrow 0\) 导出。两端的模由归纳假设为零，因为 N/JN 与 JN 都被 \(J^{m-1}\) 零化。于是有 \(\operatorname{Ext}_{\Lambda}^{i}(N,M) = 0\) ，此即所要证明的。

推论 1.--- 设 m 为整数 \textgreater0 ，J' 为 A 的包含 J\^{}m 的理想。则 \(\operatorname{prof}_{\mathrm{A}}(\mathrm{J};\mathrm{M}) \leqslant \operatorname{prof}_{\mathrm{A}}(\mathrm{J}';\mathrm{M})\) 。

事实上，\(J^{m}\) 零化 A-模 \(A/J'\) ，故 \(\operatorname{Ext}_{A}^{i}(A/J', M)\) 对每个整数 \(i < \operatorname{prof}_{A}(J; M)\) 均为零（命题 2）。

推论 2.--- 设理想 J 有限生成，并设 J' 为 A 的理想且满足 V(J) ⊃ V(J')。

\begin{enumerate}
\def\labelenumi{\alph{enumi})}
\item
  有 \(a \operatorname{prof}_{\mathrm{A}}(\mathrm{J}; \mathrm{M}) \leqslant \operatorname{prof}_{\mathrm{A}}(\mathrm{J}'; \mathrm{M})\) 。
\item
  若理想 J' 有限生成且 V(J) = V(J')，则有 \(\text{prof}_{A}(J; M) = \text{prof}_{A}(J'; M)\) 。
\end{enumerate}

根据 II, § 4, 第 3 节, 命题 11 的推论 2 及 § 2, 第 6 节, 命题 15，存在整数 m \textgreater{} 0 使 \(J^{m} \subset J'\) 。于是断言 a) 由推论 1 得出，而断言 b) 又由它导出。

当理想 J 不是有限生成时，推论 2 可能不成立（习题 2）。

\subsection{深度与零调性}

命题 3.--- 设 A 为环，C 为 A-模的下有界复形，p 为整数。假设对每对整数 \((m,n)\) （满足 \(m \geqslant n \geqslant p\) ），A-模 \(C_{m}\) 关于 \(\mathrm{H}_{n}(\mathrm{C})\) 的零化子的深度都 \textgreater m - n。则 \(\mathrm{H}_{n}(\mathrm{C}) = 0\) 对 \(n \geqslant p\) 成立。

由于 C 左侧有界，当 n 充分大时 \(\mathrm{H}_{n}(\mathrm{C})\) 为零。若结论不真，则存在整数 \(q \geqslant p\) ，使得 \(\mathrm{H}_{n}(\mathrm{C}) = 0\) 对 n \textgreater{} q 成立，而 \(\mathrm{H}_{q}(\mathrm{C}) \neq 0\) 。以 J 表示 \(\mathrm{H}_{q}(\mathrm{C})\) 的零化子；则有 \(\operatorname{prof}_{\mathrm{A}}(\mathrm{J}; \mathrm{H}_{q}(\mathrm{C})) = 0\) 。此外，由于 \(\mathrm{Z}_{q}(\mathrm{C})\) 是 \(C_{q}\) 的子模，且由假设 \(\operatorname{prof}_{\mathrm{A}}(\mathrm{J}; \mathrm{C}_{q}) > q - q = 0\) ，故有 \(\operatorname{prof}_{\mathrm{A}}(\mathrm{J}; \mathrm{Z}_{q}(\mathrm{C})) > 0\) 。由正合列

\[
0 \to \mathrm{B} _ {q} (\mathrm{C}) \to \mathrm{Z} _ {q} (\mathrm{C}) \to \mathrm{H} _ {q} (\mathrm{C}) \to 0
\]

导出等式 \(\mathrm{prof}_{\mathrm{A}}(\mathrm{J};\mathrm{B}_q(\mathrm{C})) = 1\) （第 1 节, 命题 1）。按 \(q\) 的定义，\(\mathrm{B}_n(\mathrm{C})\) 等于 \(\mathrm{Z}_n(\mathrm{C})\) ，对每个整数 \(n > q\) 皆然。由典范正合列

\[
0 \to \mathrm{B} _ {n} (\mathrm{C}) \to \mathrm{C} _ {n} \to \mathrm{B} _ {n - 1} (\mathrm{C}) \to 0 \quad (n > q)
\]

以及假设 \(\operatorname{prof}_{\mathrm{A}}(\mathrm{J};\mathrm{C}_{n}) > n - q\) ，用归纳法得到等式 \(\operatorname{prof}_{\mathrm{A}}(\mathrm{J};\mathrm{B}_{n}(\mathrm{C})) = n - q + 1\) 对每个 \(n \geqslant q\) 成立（同上引文）。但这是荒谬的，因为 \(\mathrm{B}_{n}(\mathrm{C})\) 当 \(n\) 充分大时为零。

推论 1.--- 设 A 为环，J 为 A 的理想，C 为 A-模的下有界复形，p 为整数。假设 \(\mathrm{JH}_{m}(\mathrm{C}) = 0\) 且 \(\operatorname{prof}_{\mathrm{A}}(\mathrm{J}; \mathrm{C}_{m}) > m - p\) 对 \(m \geqslant p\) 成立。则 \(\mathrm{H}_{n}(\mathrm{C}) = 0\) 对 \(n \geqslant p\) 成立。

事实上，对 \(n \geqslant p\) ，零化子 \(J_n\) （即 \(H_n(C)\) 的零化子）包含 \(J\) ；因此有 \(\mathrm{prof}_{\mathrm{A}}(\mathrm{J}_{n};\mathrm{C}_{m}) \geqslant \mathrm{prof}_{\mathrm{A}}(\mathrm{J};\mathrm{C}_{m})\) （ \(n^{\circ}\) 1, 命题 2 的推论 1），从而命题的假设得到满足。

推论 2.-- 设 A 为局部环，C 为 A-模的下有界复形，p 为整数。假设对 \(m \geqslant p\) ，\(\mathrm{H}_{m}(\mathrm{C})\) 的长度有限且 \(\mathrm{C}_{m}\) 的深度 \(>m-p\) 。则 \(\mathrm{H}_{n}(\mathrm{C}) = 0\) 对 \(n \geqslant p\) 成立。

A-模 \(\bigoplus_{m\geqslant p}\mathrm{H}_{m}(\mathrm{C})\) 的长度有限。以 J 表示其零化子；由 A, VIII, § 1, 第 3 节, 推论，环 A/J 是 Artin 环，故 J 包含 A 的极大理想的一个幂 (A, VIII, § 10, 第 1 节, 定理 1)。于是 \(\mathrm{prof}_A(\mathrm{J};\mathrm{C}_m)\geqslant \mathrm{prof}(\mathrm{C}_m) > m - p\) 对 \(m\geqslant p\) 成立（第 1 节, 命题 2 的推论 1），从而可以应用推论 1。

\subsection{深度与 Koszul 复形}

设 A 为环，M 为 A-模，\(\mathbf{x} = (x_i)_{i \in I}\) 为 A 中元素的一个族。用 \(u : \mathrm{A}^{(1)} \to \mathrm{A}\) 表示使 \(u(c_i) = x_i\) 对每个 \(i \in I\) 成立的线性形式，并以 \(\mathbf{K}^{\bullet}(\mathbf{x}, \mathrm{M})\) 表示与 u 相伴的复形 \(\mathbf{K}_{\mathrm{A}}^{\bullet}(u, \mathrm{M})\) (A, X, 第 147 页)。\(\mathbf{K}^p(\mathbf{x}, \mathrm{M}) = 0\) 对 p \textless{} 0 成立；对 \(p \geqslant 0\) ，A-模 \(\mathbf{K}^p(\mathbf{x}, \mathrm{M}) = \operatorname{Hom}_{\Lambda} (\boldsymbol{\Lambda}^p(\mathrm{A}^{(1)}), \mathrm{M})\) 典范地等同于 A-模 \(C_I^p(\mathrm{M})\) ，后者由 \(I^p\) 到 M 的交错映射构成 (A, X, 第 153 页)，而微分 \(\partial^p : \mathbf{K}^p(\mathbf{x}, \mathrm{M}) \longrightarrow \mathbf{K}^{p+1}(\mathbf{x}, \mathrm{M})\) 由下述公式给出

\[
\left(\partial^ {p} \boldsymbol {m}\right) \left(\alpha_ {1}, \dots , \alpha_ {p + 1}\right) = \sum_ {j = 1} ^ {p + 1} (- 1) ^ {j + 1} x _ {\alpha_ {j}} \boldsymbol {m} \left(\alpha_ {1}, \dots , \alpha_ {j - 1}, \alpha_ {j + 1}, \dots , \alpha_ {p + 1}\right)
\]

对 \(\boldsymbol{m} \in \mathbf{K}^{p}(\mathbf{x}, \mathrm{M})\) 及 \((\alpha_{1}, \ldots, \alpha_{p+1}) \in \mathrm{I}^{p+1}\) 成立 (A, X, 第 154 页, 公式 (12))。特别地，由此得出复形 \(\mathbf{K}^{\bullet}(\mathbf{x}, \mathrm{M})\) 仅依赖于 M 作为 Z-模的结构以及自同态 \((x_{i})_{\mathrm{M}}\) 。

以 \(\mathrm{H}^{\bullet}(\mathbf{x},\mathrm{M})\) 表示复形 \(\mathbf{K}^{\bullet}(\mathbf{x},\mathrm{M})\) 的同调。A-模 \(\mathrm{H}^{0}(\mathbf{x},\mathrm{M})\) 等同于 \(\mathrm{Hom}_{\Lambda}(\mathrm{A}/\mathrm{J},\mathrm{M})\) ，其中 J 为由 \(x_{i}\) 生成的 A 的理想 (A, X, 第 147 页, 引理 1)。

设 \((\mathrm{M}_{\alpha})_{\alpha\in\mathrm{K}}\) 为一族 A-模，M 为其乘积；复形 \(\mathbf{K}^{\bullet}(\mathbf{x},\mathrm{M})\) 典范同构于诸 \(\mathbf{K}^{\bullet}(\mathbf{x},\mathrm{M}_{\alpha})\) 的乘积复形，因而对每个整数 s，A-模 \(\mathrm{H}^{s}(\mathbf{x},\mathrm{M})\) 等同于诸 \(\mathrm{H}^{s}(\mathbf{x},\mathrm{M}_{\alpha})\) 的乘积 (A, X, 第 28 页, 命题 1)。

定理 1.-- 设 A 为环，J 为 A 的理想，\(\mathbf{x} = (x_i)_{i \in I}\) 为 J 的一个生成族，M 为 A-模。M 关于 J 的深度是 \(\mathbf{N} \cup \{+\infty\}\) 中由满足下述条件的整数 \(n\) 组成的集合的下确界：\(\Pi^n(\mathbf{x}, \mathbf{M}) \neq 0\) 。

令 \(p = \text{prof}_{\Lambda}(\text{J}; \text{M})\) 。考察复形 \(\mathbf{K}^{\bullet}(\mathbf{x}, \text{M})\) 。它的同调被 J 零化 ( \(\Lambda\) , X, 第 148 页, 推论 2)，且诸模 \(\mathbf{K}^{i}(\mathbf{x}, \text{M})\) 中每一个关于 J 的深度等于 \(p\) 或 \(+\infty\) （第 1 节, 注 4）。于是由第 2 节的推论 1 得出：\(\text{H}^{i}(\text{x}, \text{M}) = 0\) 对 \(i < p\) 成立。剩下要证明 \(\text{H}^{p}(\text{x}, \text{M})\) 在 \(p < +\infty\) 时非零。

情形 \(p = 0\) 是显然的，故设 \(0 < p < +\infty\) 且 \(\mathrm{H}^p(\mathbf{x},\mathrm{M}) = 0\) 。设 L 为 A-模 \(\mathrm{A / J}\) 的一个自由分解；以 C 表示复形 \(\mathrm{Homgr}_{\mathrm{A}}(\mathrm{L},\mathrm{M})\) 。\(\Lambda\) -模 \(\mathrm{H}^i(\mathrm{C})\) 同构于 \(\mathrm{Ext}_{\mathrm{A}}^i(\mathrm{A / J},\mathrm{M})\) (A, X, 第 100 页, 定理 1)；因而它当 i \textless{} p 时为零。~于是，对 i \textless{} p 有典范正合列

\[
0 \rightarrow \mathrm{B} ^ {i} (\mathrm{C}) \rightarrow \mathrm{C} ^ {i} \rightarrow \mathrm{B} ^ {i + 1} (\mathrm{C}) \rightarrow 0.
\]

A-模 \(C^{i}\) 是一些同构于 M 的 A-模的乘积；故 \(\mathrm{H}^{s}(\mathbf{x},\mathrm{C}^{i})=0\) 对 \(s\leqslant p\) 成立。由上述正合列及 A, X, 第 150 页 推出：连接同态 \(\partial^{s}:\mathrm{H}^{s}(\mathbf{x},\mathrm{B}^{i+1}(\mathrm{C}))\longrightarrow\mathrm{H}^{s+1}(\mathbf{x},\mathrm{B}^{i}(\mathrm{C}))\) 对 \(s\leqslant p\) 及 i\textless p 是单射；由于 \(\mathrm{B}^{0}(\mathrm{C})=0\) ，由此得出 \(\mathrm{H}^{p-i}(\mathbf{x},\mathrm{B}^{i+1}(\mathrm{C}))\) 当 i\textless p 时为零。~特别地，有 \(\mathrm{II}^{1}(\mathbf{x},\mathrm{B}^{p}(\mathrm{C}))=0\) ，于是正合列

\[
0 \to \mathrm{B} ^ {p} (\mathrm{C}) \to \mathrm{Z} ^ {p} (\mathrm{C}) \to \mathrm{H} ^ {p} (\mathrm{C}) \to 0
\]

给出一个满射 \(\mathrm{H}^{0}(\mathbf{x},\mathrm{Z}^{p}(\mathrm{C}))\longrightarrow\mathrm{H}^{0}(\mathbf{x},\mathrm{H}^{p}(\mathrm{C}))\) 。由于 \(\mathrm{H}^{p}(\mathrm{C})\) 同构于 \(\operatorname{Ext}_{\mathrm{A}}^{p}(\mathrm{A}/\mathrm{J},\mathrm{M})\) ，而后者非零且被 J 零化，故有 \(\mathrm{H}^{0}(\mathbf{x},\mathrm{H}^{p}(\mathrm{C}))\neq0\) ，从而 \(\mathrm{H}^{0}(\mathbf{x},\mathrm{Z}^{p}(\mathrm{C}))\neq0\) ，于是 \(\mathrm{H}^{0}(\mathbf{x},\mathrm{C}^{p})\neq0\) 。但这蕴含 \(\mathrm{H}^{0}(\mathbf{x},\mathrm{M})\neq0\) ，与假设矛盾。因此 \(\mathrm{H}^{p}(\mathbf{x},\mathrm{M})\neq0\) ，证明完毕。

推论 1.--- 设理想 J 有限生成且 JM ≠ M。则 prof \(_{A}\) (J; M) 有限且 ≤ Card(I)；为使它等于 Card(I)，必须且只需族 x 对 M 是完全割的 (A, X, 第 157 页, 定义 2)。

首先设 I 有限，并以 \(r\) 表示其基数。A-模 \(\mathrm{H}^r (\mathbf{x},\mathrm{M})\) 典范地同构于 \(\mathrm{H}_0(\mathbf{x},\mathrm{M})\) ，后者又同构于 M/JM (A, X, 第 155 页)；因而不等式 \(\mathrm{prof}_{\Lambda}(\mathrm{J};\mathrm{M})\leqslant r\) 由定理 1 得出。为使等号成立，必须且只需 A-模 \(\mathrm{H}^i (\mathbf{x},\mathrm{M})\) 对 \(i < r\) 均为零，而这恰好意味着族 \(\mathbf{x}\) 对 M 是完全割的 (A, X, 第 157 页)。

根据前文，\(\mathrm{prof}_{\mathrm{A}}(\mathrm{J};\mathrm{M})\) 有限；剩下要证明：若族 \(\mathbf{x}\) 对 M 完全割，则集 I 有限。但条件 \(\mathrm{H}_{1}(\mathbf{x},\mathrm{M}) = 0\) (A, X, 第 157 页, 定义 2) 蕴含我们有正合列

\[
\boldsymbol {\Lambda} ^ {2} (\mathrm{A} ^ {(\mathrm{I})}) \otimes_ {\mathrm{A}} \mathrm{M} \xrightarrow {\partial_ {2}} \mathrm{M} ^ {(\mathrm{I})} \xrightarrow {\partial_ {1}} \mathrm{JM} \to 0,
\]

其中 \(\partial_{2}\) 的像含于 \(\mathrm{JM}^{(1)}\) 。用 A/J 作张量积，我们导出从 \((\mathrm{M} / \mathrm{JM})^{(1)}\) 到 \(\mathrm{JM} / \mathrm{J}^2\mathrm{M}\) 上的一个 A-线性同构。而后一模是有限型的，因为 J 与 M 都是；由于 M/JM 非零，于是得出集 I 有限。

推论 2. 设 A 为局部环，J 为 A 的异于 A 的有限生成理想，M 为非零的有限生成 A-模。令 \(r = [\mathrm{J} / \mathfrak{m}_{\mathrm{A}}\mathrm{J} : \kappa_{\mathrm{A}}]\) 。有 \(\operatorname{prof}_{\Lambda}(\mathrm{J};\mathrm{M}) \leqslant r\) ；当且仅当 J 由一个对 M 完全割的族生成时等号成立。这时，为使 J 的一个生成族是完全割的，必须且只需它含 r 个元素。

由 Nakayama 引理，有 \(\mathrm{JM} \neq \mathrm{M}\) ，且 \(r\) 是 \(\mathbf{J}\) 的生成元的最小个数；因此推论 2 由推论 1 得出。

命题 4.--- 设 \(\rho: A \to B\) 为环同态，\(J\) 为 \(A\) 的理想，\(N\) 为 B-模。我们有等式 \(\operatorname{prof}_{\mathrm{A}}(\mathrm{J}; \mathrm{N}) = \operatorname{prof}_{\mathrm{B}}(\mathrm{JB}; \mathrm{N})\) 。

设 \(\mathbf{x} = (x_{i})_{i \in \mathbb{I}}\) 为 J 的一个生成族；族 \(\rho(\mathbf{x}) = (\rho(x_{i}))_{i \in \mathbb{I}}\) 生成 JB。按构造，复形 \(\mathbf{K}^{\bullet}(\boldsymbol{\rho}(\mathbf{x}), \mathbf{N})\) 等于 \(\mathbf{K}^{\bullet}(\mathbf{x}, \mathbf{N})\) 。因此本命题由定理 1 得出。

推论. 设 A 为局部环，a 为 A 的异于 A 的理想，M 为被 a 零化的 A-模。我们有 \(\operatorname{prof}_{\mathrm{A}}(\mathrm{M}) = \operatorname{prof}_{\mathrm{A}/\mathfrak{a}}(\mathrm{M})\) 。

设 \(\rho: A \to B\) 为环同态，\(x = (x_i)_{i \in I}\) 为 \(\Lambda\) 中元素的一个有限族，M 为 A-模。对每个整数 \(p\) ，以 \(u^p: B \otimes_\Lambda C_I^p(M) \longrightarrow C_I^p(B \otimes_\Lambda M)\) 表示把 \(b \otimes m\) 对应到交错映射 \((\alpha_1, \ldots, \alpha_p) \mapsto b \otimes m(\alpha_1, \ldots, \alpha_p)\) 的 B-线性同态。族 \((u^p)\) 定义了复形之间的一个同构

\[
u: \mathrm{B} \otimes_ {\mathrm{A}} \mathbf {K} ^ {\bullet} (\mathbf {x}, \mathrm{M}) \longrightarrow \mathbf {K} ^ {\bullet} (\mathbf {x}, \mathrm{B} \otimes_ {\mathrm{A}} \mathrm{M}).
\]

考察典范同态

\[
\gamma^ {p} (\mathrm{B}, \mathbf {K} ^ {\bullet} (\mathbf {x}, \mathrm{M})): \mathrm{B} \otimes_ {\mathrm{A}} \mathrm{H} ^ {p} (\mathbf {x}, \mathrm{M}) \longrightarrow \mathrm{H} ^ {p} (\mathrm{B} \otimes_ {\mathrm{A}} \mathbf {K} ^ {\bullet} (\mathbf {x}, \mathrm{M}))
\]

(A, X, 第 62 页)；与 \(\mathrm{H}^{p}(u)\) 复合，由它导出同态

\[
v ^ {p}: \mathrm{B} \otimes_ {\Lambda} \mathrm{H} ^ {p} (\mathbf {x}, \mathrm{M}) \longrightarrow \mathrm{H} ^ {p} (\mathbf {x}, \mathrm{B} \otimes_ {\Lambda} \mathrm{M}).
\]

引理 1.-- 若 A-模 B 平坦，则同态 \(v^p\) 对每个整数 \(p\) 都是双射。

这由 A, X, 第 66 页, 推论 2 得出。

命题 5.--- 设 A 为环，J 为 A 的有限生成理想，M 为 A-模。以 Ω 表示属于 Supp(M) 且包含 J 的 A 的极大理想所成的集。我们有

\[
\operatorname{prof} _ {A} (J; M) = \inf _ {\mathfrak {p} \in \operatorname{Spec} (A)} \operatorname{prof} _ {A _ {\mathfrak {p}}} \left(J _ {\mathfrak {p}}; M _ {\mathfrak {p}}\right) = \inf _ {m \in \Omega} \operatorname{prof} _ {\Lambda_ {m}} \left(J _ {m}; M _ {m}\right).
\]

设 \(\mathbf{x} = (x_{i})_{i \in 1}\) 为 J 的一个有限生成族。设 p 为 A 的素理想；理想 \(J_{p}\) 由像 \(x_{p}\) ——即族 \((x_{i})\) 在 \(A_{p}\) 中的像——生成。对每个 \(p \geqslant 0\) ，\(A_{p}\) -模 \(\left(\mathrm{H}^{p}(\mathbf{x},\mathrm{M})\right)_{\mathfrak{p}}\) 同构于 \(\mathrm{H}^{p}(\mathbf{x}_{\mathfrak{p}},\mathrm{M}_{\mathfrak{p}})\) （引理 1）；因此由定理 1，有 \(\operatorname{prof}_{\Lambda}(\mathrm{J};\mathrm{M}) \leqslant \inf_{\mathfrak{p} \in \operatorname{Spec}(\mathrm{A})} \operatorname{prof}_{\mathrm{A}_{\mathfrak{p}}}(\mathrm{J}_{\mathfrak{p}}; \mathrm{M}_{\mathfrak{p}}) \leqslant \inf_{\mathfrak{m} \in \Omega} \operatorname{prof}_{\mathrm{A}_{\mathfrak{m}}}(\mathrm{J}_{\mathfrak{m}}; \mathrm{M}_{\mathfrak{m}})\) 。

设 \(p\) 为一个整数，它严格小于 \(\operatorname{prof}_{\mathrm{A}_{\mathfrak{m}}}\left(\mathrm{JA}_{\mathfrak{m}};\mathrm{M}_{\mathfrak{m}}\right)\) （对每个 \(\mathfrak{m}\in\Omega\) ）。则 \(\mathrm{H}^{p}\left(\mathbf{x}_{\mathfrak{m}},\mathrm{M}_{\mathfrak{m}}\right)=0\) 对 A 的每个极大理想 \(\mathfrak{m}\) 成立：当 \(\mathfrak{m}\in\Omega\) 时这由定理 1 得出，由 \(\mathrm{M}_{\mathfrak{m}}=0\) （当 \(\mathfrak{m}\notin\operatorname{Supp}(\mathrm{M})\) 时）这一事实亦可得之，而由下述事实亦可得之：理想 \(\mathrm{JA}_{\mathfrak{m}}\) 零化 \(\mathrm{H}^{p}\left(\mathbf{x}_{\mathfrak{m}},\mathrm{M}_{\mathfrak{m}}\right)\) (A, X, 第 148 页, 推论 2)，且等于 \(\mathrm{A}_{\mathfrak{m}}\) 当 \(\mathfrak{m}\notin\mathrm{V}(\mathrm{J})\) 时。于是 \(\left(\mathrm{H}^{p}(\mathbf{x},\mathrm{M})\right)_{\mathfrak{m}}=0\) 对 A 的每个极大理想 \(\mathfrak{m}\) 成立，这蕴含 \(\mathrm{H}^{p}(\mathbf{x},\mathrm{M})=0\) (II, § 3, 第 3 节, 定理 1 的推论 2)。本命题于是由定理 1 得出。

命题 6.--- 设 A 为环，J 为 A 的有限生成理想，M 为 A-模。设 B 为环，\(\rho: A \to B\) 为使 B 成为平坦 A-模的环同态。

\begin{enumerate}
\def\labelenumi{\alph{enumi})}
\item
  有 \(a \operatorname{prof}_{\mathrm{A}}(\mathrm{J}; \mathrm{M}) \leqslant \operatorname{prof}_{\mathrm{B}}(\mathrm{JB}; \mathrm{B} \otimes_{\mathrm{A}} \mathrm{M})\) 。
\item
  再设 Supp(M) 的每个包含 J 的极大理想都属于典范映射 \(\operatorname{Spec}(B) \to \operatorname{Spec}(A)\) 的像。则 \(\operatorname{prof}_{\mathrm{A}}(\mathrm{J}; \mathrm{M}) = \operatorname{prof}_{\mathrm{B}}(\mathrm{JB}; \mathrm{B} \otimes_{\mathrm{A}} \mathrm{M})\) 。特别地，当 A-模 B 忠实平坦时即属此情形。
\end{enumerate}

断言 a) 由定理 1 与引理 1 得出。

设 \(p\) 为严格小于 \(\text{prof}_B(\text{JB};\text{B} \otimes_A \text{M})\) 的整数，并设 \(\mathfrak{m}\) 为属于 \(\text{Supp(M)} \cap \text{V(J)}\) 的 A 的极大理想。设 \(\mathbf{x}\) 为理想 J 的一个有限生成族。在 b) 的假设下，存在 B 的素理想 \(\mathfrak{n}\) 位于 \(\mathfrak{m}\) 之上，并且有一个典范同构

\[
\mathrm{B} _ {\mathfrak {n}} \otimes_ {\mathrm{A} _ {\mathfrak {m}}} \left(\mathrm{A} _ {\mathfrak {m}} \otimes_ {\Lambda} \mathrm{H} ^ {p} (\mathbf {x}, \mathrm{M})\right) \longrightarrow \mathrm{B} _ {\mathfrak {n}} \otimes_ {\mathrm{B}} \left(\mathrm{B} \otimes_ {\mathrm{A}} \mathrm{H} ^ {p} (\mathbf {x}, \mathrm{M})\right).
\]

现在 \(\mathrm{B}\otimes_{\Lambda}\mathrm{H}^{p}(\mathbf{x},\mathrm{M})\) 同构于 \(\mathrm{H}^{p}(\mathsf{p}(\mathbf{x}),\mathrm{B}\otimes_{\Lambda}\mathrm{M})\) （引理 1），故为零；此外 \(B_{n}\) 在 \(A_{m}\) 上忠实平坦 (I, § 3, 第 5 节, 命题 9 及 II, § 3, 第 4 节, 命题 14 及注)。于是有 \(A_{m}\otimes_{\Lambda}\mathrm{H}^{p}(\mathbf{x},\mathrm{M})=0\) ，从而 p \textless{} prof \(_{A_{m}}(J_{m};M_{m})\) （引理 1 与定理 1）。b) 的第一个断言于是由命题 5 得出；第二个断言由 I, § 3, 第 5 节, 命题 9 得出）。

推论.--- 设 A 为 Noether 环，J 为 A 的理想，M 为有限生成 A-模，\(\widehat{\Lambda}\) 与 \(\widehat{M}\) 为 A 与 M 关于 J-进拓扑的分离完备化。则有 \(\operatorname{prof}_{\mathbf{A}}(\mathbf{J};\mathbf{M}) = \operatorname{prof}_{\widehat{\mathbf{A}}}(\widehat{\mathbf{J}}\widehat{\mathbf{A}};\widehat{\mathbf{M}})\) 。

事实上，A-模 \(\widehat{\mathsf{A}}\) 平坦，且 \(\widehat{\mathsf{A}}\) -模 \(\widehat{\mathsf{M}}\) 同构于 \(\widehat{\mathsf{A}}\otimes_{\mathsf{A}}\mathsf{M}\) (III, § 3, 第 4 节, 定理 3)；此外，A 的每个包含 J 的极大理想都属于映射 \(\operatorname{Spec}(\widehat{\mathsf{A}})\to\operatorname{Spec}(\mathsf{A})\) 的像（同上引文, 命题 8）。

\subsection{深度与正则序列}

设 A 为环，M 为 A-模。回忆（A, X, 第 158 页），A 中元素的有限序列 \((x_{1},\ldots,x_{r})\) 称为关于 M 正则的，或 M-正则的，如果对 \(i=1,\ldots,r\) ，乘以 \(x_{i}\) 的映射在 A-模 \(\mathrm{M}/(x_{1}\mathrm{M}+\ldots+x_{i-1}\mathrm{M})\) 中是单射。设 \((x_{1},\ldots,x_{r})\) 为 M-正则序列；对每个平坦 A-模 N，序列 \((x_{1},\ldots,x_{r})\) 是 \(M\otimes_{A}N\) -正则的。若 \(\rho:A\to B\) 是使 B 成为平坦 A-模的环同态，则序列 \((\rho(x_{1}),\ldots,\rho(x_{r}))\) 对 B-模 \(B\otimes_{A}M\) 正则。特别地，对 A 的每个素理想 p，在 \(A_{p}\) 中序列 \((x_{1},\ldots,x_{r})\) 的像是 \(M_{p}\) -正则的。

在下文中我们主要考虑 M-正则序列这一概念在下列情形的形态：环 A 为局部 Noether 环，模 M 有限生成，且序列的元素属于 \(m_{A}\) ；此时 M-正则序列的概念与 M 的完全割序列的概念一致 (A, X, 第 160 页, 推论 1)。

命题 7.--- 设 A 为环，J 为 A 的理想，M 为 A-模，\((x_{1},\ldots,x_{r})\) 为由 J 中元素组成的 M-正则序列。则有

\[
\operatorname{prof} _ {\mathrm{A}} (\mathrm{J}; \mathrm{M}) = r + \operatorname{prof} _ {\mathrm{A}} (\mathrm{J}; \mathrm{M} / (x _ {1} \mathrm{M} + \dots + x _ {r} \mathrm{M}))
\]

特别地 \(\operatorname{prof}_{\mathrm{A}}(\mathrm{J};\mathrm{M})\geqslant r\)

\(r=1\) 的情形由第 1 节的注 5 应用于正合序列而得出

\[
0 \rightarrow \mathrm{M} \xrightarrow {(x _ {1}) _ {\mathrm{M}}} \mathrm{M} \longrightarrow \mathrm{M} / x _ {1} \mathrm{M} \rightarrow 0.
\]

一般情形由对 r 作归纳法从它推出。

定理 2.--- 设 A 为 Noether 环，J 为 A 的理想，M 为有限生成 A-模。

\begin{enumerate}
\def\labelenumi{\alph{enumi})}
\item
  设 \(\mathrm{prof}_{\mathrm{A}}(\mathrm{J};\mathrm{M})\) 有限。则每个由 J 中元素组成的 M-正则序列都可延拓为一个由 J 中元素组成、长度为 \(\mathrm{prof}_{\mathrm{A}}(\mathrm{J};\mathrm{M})\) 的 M-正则序列。
\item
  M 关于 J 的深度是由 J 中元素组成的 M-正则序列的长度的上确界。
\item
  \(\operatorname{prof}_{\mathrm{A}}(\mathrm{J};\mathrm{M})\) 有限必须且只需 M 的支撑集与 \(\mathrm{V}(\mathrm{J})\) 相交，或等价地，\(\mathrm{JM} \neq \mathrm{M}\) 。
\end{enumerate}

设 \((x_{1},\ldots,x_{r})\) 为由 J 中元素组成的 M-正则序列。有 \(r\leqslant\operatorname{prof}_{\mathrm{A}}(\mathrm{J};\mathrm{M})\) （命题 7）；设该不等式是严格的。用 N 表示 A-模 \(\mathrm{M}/(x_{1}\mathrm{M}+\ldots+x_{r}\mathrm{M})\) 。有 \(\operatorname{prof}_{\mathrm{A}}(\mathrm{J};\mathrm{N})>0\) （同上引文），故存在 J 的元素 x 使乘以 x 的映射 \(x_{N}\) 是单射（第 1 节, 注 2），即序列 \((x_{1},\ldots,x_{r},x)\) 是 M-正则的。由归纳法得出：对每个满足 \(r\leqslant s\leqslant\operatorname{prof}_{\mathrm{A}}(\mathrm{J};\mathrm{M})\) 的整数 s，序列 \((x_{1},\ldots,x_{r})\) 可补全为长度 s 的 M-正则序列，这就蕴含断言 a) 与 b)。断言 c) 由第 1 节的注 1 及第 3 节的定理 1 的推论 1 得出。

推论 1. 对每个由 J 中元素组成的 M-正则序列 \((x_{1},\ldots ,x_{r})\) ，下列性质等价：

\begin{enumerate}
\def\labelenumi{(\roman{enumi})}
\item
  有 \(r = \mathrm{prof}_{\mathrm{A}}(\mathrm{J};\mathrm{M})\)
\item
  序列 \((x_{1},\ldots,x_{r})\) 在由 J 中元素组成的 M-正则序列之中是极大的；
\item
  A-模 \(\mathrm{M}/(x_{1}\mathrm{M}+\ldots+x_{r}\mathrm{M})\) 有一个被 J 零化的非零元素；
\end{enumerate}

\[
\text {(iv)} \text {   we have   } \operatorname{Ass} \bigl (\mathrm{M} / (x _ {1} \mathrm{M} + \ldots + x _ {r} \mathrm{M}) \bigr) \cap \mathrm{V} (\mathrm{J}) \neq \emptyset .
\]

 与 的等价性由定理 2 得出；、 与 的等价性由第 1 节的注 2 应用于 A-模 M/(x₁M + \ldots{} + xᵣM) 而得出。

推论 2. --- 设 A 为 Noether 局部环，M 为非零的有限生成 A-模。则

\[
\operatorname{prof} _ {\mathrm{A}} (\mathrm{M}) \leqslant \dim_ {\mathrm{A}} (\mathrm{M}) <   + \infty .
\]

事实上，每个由 \(\mathfrak{m}_{\mathrm{A}}\) 中元素组成的 M-正则序列对 M 是完全割的 (A, X, 第 157 页, 命题 5)，从而对 M 是割的 (VIII, § 3, 第 2 节, 命题 3 的推论)；因此其长度以整数 \(\dim_{\Lambda}(\mathrm{M})\) 为上界（同上引文, 定理 1）。

\subsection{沿闭子集的深度}

设 A 为 Noether 环，F 为 \(\operatorname{Spec}(A)\) 的闭子集，M 为 A-模。由第 1 节的命题 2 的推论 2，元素 \(\operatorname{prof}_{A}(J;M)\) （属于 \(\mathbf{N} \cup \{+\infty\}\) ）不依赖于满足 \(F = V(J)\) 的 A 的理想 J；它称为 M 沿 F 的深度，记作 \(\operatorname{prof}_{F}(M)\) 。

注.-- 1) 设 \(r\) 为整数。由第 1 节的命题 2 及 II, § 4, 第 4 节, 命题 17 的推论 2，不等式 \(\mathrm{prof}_{\mathrm{F}}(\mathbf{M}) \geqslant r\) 等价于下述性质：对每个支撑集包含于 F 的有限生成 A-模 N，\(\mathrm{Ext}_{\mathrm{A}}^{i}(\mathrm{N},\mathrm{M}) = 0\) 对 \(i < r\) 成立。

\begin{enumerate}
\def\labelenumi{\arabic{enumi})}
\setcounter{enumi}{1}
\tightlist
\item
  设 A-模 M 有限生成。由第 1 节的注 2 与第 4 节的定理 2，我们有下列等价关系
\end{enumerate}

\[
\begin{array}{c} \operatorname{prof} _ {\mathrm{F}} (\mathrm{M}) = 0 \Longleftrightarrow \operatorname{Ass} (\mathrm{M}) \cap \mathrm{F} \neq \emptyset \\ \operatorname{prof} _ {\mathrm{F}} (\mathrm{M}) <   + \infty \Longleftrightarrow \operatorname{Supp} (\mathrm{M}) \cap \mathrm{F} \neq \emptyset . \end{array}
\]

命题 8. 设 A 为 Noether 环，F 为 \(\operatorname{Spec}(\Lambda)\) 的闭子集，M 为有限生成 A-模。我们有

\[
\operatorname{prof} _ {F} (M) = \inf _ {\mathfrak {p} \in F} \operatorname{prof} _ {A _ {\mathfrak {p}}} \left(M _ {\mathfrak {p}}\right) = \inf _ {\mathfrak {p} \in \operatorname{Supp} (M) \cap F} \operatorname{prof} _ {A _ {\mathfrak {p}}} \left(M _ {\mathfrak {p}}\right).
\]

若 \(\operatorname{prof}_{F}(M) = +\infty\) 这是显然的（注 2）。若 \(\operatorname{prof}_{F}(M) = 0\) ，则存在素理想 \(\mathfrak{p} \in \operatorname{Ass}(M) \cap F\) （注 2）；我们有 \(\mathfrak{p}A_{\mathfrak{p}} \in \operatorname{Ass}(M_{\mathfrak{p}})\) (IV, § 1, 第 2 节, 命题 5)，从而 \(\operatorname{prof}_{\Lambda_{\mathfrak{p}}}(\mathrm{M}_{\mathfrak{p}}) = 0\) （第 1 节的注 2），命题在此情形得证。

设 \(0 < \operatorname{prof}_{\mathrm{F}}(\mathrm{M}) < +\infty\) 。设 J 为 A 的理想使 \(\mathrm{V}(\mathrm{J}) = \mathrm{F}\) ，并设 \(x\) 为 J 中使乘以 x 的映射 \(x_{\mathrm{M}}\) 为单射的元素（同上引文）。对每个素理想 \(\mathfrak{p}\) ，乘以 x 的映射 \(x_{\mathrm{M}_{\mathfrak{p}}}\) 是单射的。于是由第 4 节的命题 7，我们有

\[
\begin{array}{c} \operatorname{prof} _ {\mathrm{F}} (\mathrm{M} / x \mathrm{M}) = \operatorname{prof} _ {\mathrm{F}} (\mathrm{M}) - 1 \\ \operatorname{prof} _ {\Lambda_ {\mathfrak {p}}} \big ((\mathrm{M} / x \mathrm{M}) _ {\mathfrak {p}} \big) = \operatorname{prof} _ {\Lambda_ {\mathfrak {p}}} (\mathrm{M} _ {\mathfrak {p}}) - 1. \end{array}
\]

然后对整数 \(\mathrm{prof}_{\mathbb{F}}(\mathbb{M})\) 作归纳即得结论。

注 3. 若 q 是 Supp(M) 的点，则我们有 \(\operatorname{prof}_{\mathrm{A}}(\mathfrak{q};\mathrm{M}) = \inf_{\mathfrak{p}\supseteq q} \operatorname{prof}_{\Lambda_{\mathfrak{p}}}(\mathrm{M}_{\mathfrak{p}})\) 。特别地，有不等式 \(\operatorname{prof}_{\mathrm{A}}(\mathfrak{q};\mathrm{M}) \leqslant \operatorname{prof}_{\mathrm{A}_{\mathfrak{q}}}(\mathrm{M}_{\mathfrak{q}})\) ；当 q 极大时取等号。在一般情形，可能有 \(\operatorname{prof}_{\mathrm{A}}(\mathfrak{q};\mathrm{M}) < \operatorname{prof}_{\Lambda_{\mathfrak{q}}}(\mathrm{M}_{\mathfrak{q}})\) ；也可能有 \(\operatorname{prof}_{\mathrm{A}}(\mathfrak{q};\mathrm{M}) < \inf \operatorname{prof}_{\mathrm{A}_{\mathfrak{m}}}(\mathrm{M}_{\mathfrak{m}})\) ，其中 m 遍历 A 的包含 q 的极大理想组成的集合。例如，设 p 为 A 的一个非极大素理想，它包含 q 且异于 q；令 M = A/p。则有 \(\operatorname{prof}_{\mathrm{A}}(\mathfrak{q};\mathrm{M}) = 0\) ，\(\operatorname{prof}_{\mathrm{A}_{\mathfrak{q}}}(\mathrm{M}_{\mathfrak{q}}) = +\infty\) ，且对 A 的每个极大理想 m 有 \(\operatorname{prof}_{\mathrm{A}_{\mathfrak{m}}}(\mathrm{M}_{\mathfrak{m}}) > 0\) 。

命题 9.--- 设 A 为 Noether 环，M 与 N 为两个有限生成 A-模，F 为 N 的支撑集。则 \(\operatorname{prof}_{\mathbf{F}}(\mathbf{M})\) 是 \(\mathbf{N} \cup \{+\infty\}\) 中由满足下述条件的整数 \(n\) 组成的集合的下确界：\(\operatorname{Ext}_{\Lambda}^{n}(\mathbf{N}, \mathbf{M}) \neq 0\) 。

由注 1，\(\mathrm{Ext}_{\mathrm{A}}^{i}(\mathrm{N},\mathrm{M}) = 0\) 对每个 \(i < \mathrm{prof}_{\mathrm{F}}(\mathrm{M})\) 成立。剩下要证明：若 \(\mathrm{prof}_{\mathrm{F}}(\mathrm{M}) = n < +\infty\) ，则 \(\mathrm{Ext}_{\mathrm{A}}^{n}(\mathrm{N},\mathrm{M}) \neq 0\) 。设 J 为 N 的零化子；有 \(\mathrm{F} = \mathrm{V}(\mathrm{J})\) ，从而 \(\mathrm{prof}_{\mathrm{F}}(\mathrm{M}) = \mathrm{prof}_{\mathrm{A}}(\mathrm{J};\mathrm{M})\) 。由定理 2 的推论 1（ \(n^{\circ}4\) ），存在 M-正则序列 ( \(x_{1},\ldots,x_{n}\) )，其长度为 \(n\) ，由 J 中元素组成，且 A-模 \(\overline{\mathrm{M}} = \mathrm{M}/(x_{1}\mathrm{M}+\ldots+x_{n}\mathrm{M})\) 关于 J 的深度为零。由 A, X, 第 166 页, 命题 9，只需证明 \(\mathrm{Hom}_{\mathrm{A}}(\mathrm{N},\overline{\mathrm{M}})\) 非零。而由命题 8，存在 \(\mathfrak{p} \in \mathrm{Supp}(\mathrm{M}) \cap \mathrm{Supp}(\mathrm{N})\) 使 \(\mathrm{prof}_{\mathrm{A}_{\mathfrak{p}}}(\overline{\mathrm{M}}_{\mathfrak{p}}) = 0\) ，即 \(\mathrm{Hom}_{\mathrm{A}_{\mathfrak{p}}}(\kappa(\mathfrak{p}),\overline{\mathrm{M}}_{\mathfrak{p}}) \neq 0\) 。由于 \(\mathbf{N}_{\mathfrak{p}}\) 非零，\(\kappa(\mathfrak{p})\) -向量空间 \(\mathbf{N}_{\mathfrak{p}} \otimes_{\mathbf{A}_{\mathfrak{p}}} \kappa(\mathfrak{p})\) 非零（Nakayama 引理），其对偶亦然；因此存在一个满 \(\mathbf{A}_{\mathfrak{p}}\) -线性映射，它从 \(\mathbf{N}_{\mathfrak{p}}\) 映到 \(\kappa(\mathfrak{p})\) 。由此得出 \(\mathrm{Hom}_{\mathrm{A}_{\mathfrak{p}}}(\mathbf{N}_{\mathfrak{p}},\overline{\mathbf{M}}_{\mathfrak{p}}) \neq 0\) ，从而 \(\mathrm{Hom}_{\mathrm{A}}(\mathrm{N},\overline{\mathrm{M}}) \neq 0\) (II, § 2, \(n^{\circ}7\) , 命题 19)，此即所要证明的。

注 4.--- 设 A 为 Noether 环，N 为有限生成 A-模。在 \(\mathbf{N}\cup\{+\infty\}\) 中取由整数 \(n\) （使 \(\operatorname{Ext}_{\mathrm{A}}^{n}(\mathrm{N},\mathrm{A})\) 非零）组成的集合的最大下界，人们有时把它称为 N 的级并记作 grade (N)。由命题 9，它也是 A 沿 N 的支撑集的深度，或者（ \(n^{\circ}\) 4, 定理 2）是 N 的零化子的元素组成的 A-正则序列的长度的集合的最小上界。由于对 A 的每个素理想 p，\(N_{p}\) 的零化子等于 \(\operatorname{Ann}(N)_{p}\) (II, § 2, \(n^{\circ}\) 4, 公式 (9))，由 \(n^{\circ}\) 3 的命题 5 可推出等式

\[
\operatorname{grade} (\mathrm{N}) = \inf _ {\mathfrak {p} \in \operatorname{Spec} (\Lambda)} \operatorname{grade} \left(\mathrm{N} _ {\mathfrak {p}}\right) = \inf _ {\mathfrak {m} \in \Omega} \operatorname{grade} \left(\mathrm{N} _ {\mathfrak {m}}\right),
\]

其中 Ω 表示 A 的极大理想组成的集合。

\subsection{代数的深度}

引理 2.--- 设 \(\rho: A \to B\) 为 Noether 局部环之间的局部同态，N 为有限生成 B-模，y 为 \(\mathfrak{m}_{\mathrm{B}}\) 的元素。下列两个条件等价：

\begin{enumerate}
\def\labelenumi{(\roman{enumi})}
\item
  A-模 N/yN 平坦且乘以 y 的映射 \(y_{N}\) 单射；
\item
  A-模 N 平坦且乘以 y 的映射 \(y_{\kappa_{A}\otimes N}\) 单射。
\end{enumerate}

当这些条件满足时，对每个 A-模 M，乘以 y 的映射 \(y_{\mathrm{M} \otimes_{\Lambda}\mathrm{N}}\) 都单射。

设 (i) 的假设成立，并同时证明 (ii) 及最后一个断言。设 M 为 A-模。由于 A-模 N/yN 平坦，由正合序列 \(0 \rightarrow N \xrightarrow{y_{N}} N \longrightarrow N/yN \rightarrow 0\) 推出正合序列

\[
0 \rightarrow \mathrm{M} \otimes_ {\mathrm{A}} \mathrm{N} \xrightarrow {u} \mathrm{M} \otimes_ {\mathrm{A}} \mathrm{N} \longrightarrow \mathrm{M} \otimes_ {\mathrm{A}} (\mathrm{N} / y \mathrm{N}) \rightarrow 0
\]

\[
0 \to \operatorname{Tor} _ {1} ^ {\mathrm{A}} (\mathrm{M}, \mathrm{N}) \stackrel {v} {\longrightarrow} \operatorname{Tor} _ {1} ^ {\mathrm{A}} (\mathrm{M}, \mathrm{N}) \to 0
\]

其中 \(u = 1_{\mathrm{M}} \otimes y_{\mathrm{N}}\) ， \(v = \operatorname{Tor}_1^{\mathrm{A}}(1_{\mathrm{M}}, y_{\mathrm{N}})\) ；由此得出，乘以 \(y\) 的映射在 \(\mathrm{M} \otimes_{\mathrm{A}} \mathrm{N}\) 中是单射的，且在 \(\operatorname{Tor}_1^{\mathrm{A}}(\mathrm{M}, \mathrm{N})\) 中是双射的。再设 A-模 \(\mathrm{M}\) 是有限型的；则 B-模 \(\operatorname{Tor}_1^{\mathrm{A}}(\mathrm{M}, \mathrm{N})\) 是有限型的 (A, X, 第 107 页, 命题 6)，由于 \(y\) 属于 \(\mathfrak{m}_{\mathrm{B}}\) ，它由 Nakayama 引理为零，这蕴含 A-模 \(\mathrm{N}\) 平坦 (A, X, 第 74 页, 定理 2)。

(ii)⇒(i)：设 的假设成立。考察 B-模的下列两个正合序列

\begin{enumerate}
\def\labelenumi{(\arabic{enumi})}
\tightlist
\item
\end{enumerate}

\[
0 \to \mathrm{Ker} (y _ {\mathrm{N}}) \longrightarrow \mathrm{N} \stackrel {p} {\longrightarrow} \mathrm{Im} (y _ {\mathrm{N}}) \to 0\tag{1}
\]

\[
0 \to \mathrm{Im} (y _ {\mathrm{N}}) \stackrel {i} {\longrightarrow} \mathrm{N} \longrightarrow \mathrm{N} / y \mathrm{N} \to 0,\tag{2}
\]

其中 \(p\) 与 \(i\) 是典范同态。由此得出，同态 \(1 \otimes p : \kappa_{\mathrm{A}} \otimes_{\Lambda} \mathrm{N} \longrightarrow \kappa_{\mathrm{A}} \otimes_{\mathrm{A}} \mathrm{Im}(y_{\mathrm{N}})\) 是满的，且（由于 N 平坦）同态 \(1 \otimes i : \kappa_{\mathrm{A}} \otimes_{\mathrm{A}} \mathrm{Im}(y_{\mathrm{N}}) \longrightarrow \kappa_{\mathrm{A}} \otimes_{\Lambda} \mathrm{N}\) 的核同构于 \(\operatorname{Tor}_1^{\mathrm{A}}(\kappa_{\mathrm{A}}, \mathrm{N}/y\mathrm{N})\) 。而映射 \((1 \otimes i) \circ (1 \otimes p)\) ——它等于 \(y_{\kappa_{\mathrm{A}} \otimes_{\mathrm{A}} \mathrm{N}}\) ——按假设是单射；由此得出 \(1 \otimes p\) 双射且 \(1 \otimes i\) 单射，从而有 \(\operatorname{Tor}_1^{\mathrm{A}}(\kappa_{\mathrm{A}}, \mathrm{N}/y\mathrm{N}) = 0\) 。于是得出 A-模 \(\mathrm{N}/y\mathrm{N}\) 平坦 (III, § 5, 第 2 节, 定理 1 及第 4 节, 命题 2)。

由于 N 与 N/yN 在 A 上平坦，\(\mathrm{Im}(y_{\mathrm{N}})\) 亦然（正合序列 (2)）。于是由正合序列 (1) 推出，\(\kappa_{A} \otimes_{A} \mathrm{Ker}(y_{\mathrm{N}})\) 同构于 \(1 \otimes p\) 的核；故它为零；从而乘以 y 的映射 \(y_{N}\) 由 Nakayama 引理是单射的。

命题 10. 设 \(\rho: \Lambda \to B\) 为 Noether 局部环之间的局部同态，N 为有限生成 B-模，\(y = (y_1, \ldots, y_s)\) 为 \(\mathfrak{m}_{\mathrm{B}}\) 中元素组成的序列。设 \(\mathfrak{y}\) 为由该序列生成的 B 的理想。下列条件等价：

\begin{enumerate}
\def\labelenumi{(\roman{enumi})}
\item
  A-模 N/\(\mathfrak{y}\)N 平坦且序列 \(y\) 是 N-正则的；
\item
  A-模 N 平坦且序列 \(y\) 是 \((\kappa_{A} \otimes_{A} N)\) -正则的。
\end{enumerate}

当这些条件满足时，对每个 A-模 M，序列 \(y\) 是 \(M\otimes_{\mathrm{A}}N\) -正则的。

我们对 s 作归纳来证明 与 的等价性。s = 0 的情形是显然的，设 \(s \geqslant 1\) 。用 \(y'\) 表示序列 \((y_{1}, \ldots, y_{s-1})\) ，用 \(\mathfrak{y}'\) 表示由它生成的 B 的理想。把引理 2 应用于 B-模 N/ \(y'\) N 及 B 的元素 \(y_{s}\) ，条件 等价于

(i') A-模 N/ \(\mathfrak{y}'\) N 平坦，序列 \(y'\) 是 N-正则的，且乘以 \(y_s\) 的映射在 \(\kappa_{\mathrm{A}} \otimes_{\mathrm{A}} (\mathrm{N}/\mathfrak{y}'\mathrm{N}) = (\kappa_{\mathrm{A}} \otimes_{\mathrm{A}} \mathrm{N})/\mathfrak{y}'(\kappa_{\mathrm{A}} \otimes_{\mathrm{A}} \mathrm{N})\) 中是单射的。

由归纳假设，该条件等价于。

最后一个断言类似地由引理 2 的最后一个断言对 \(s\) 作归纳得出。

命题 11. 设 \(\rho: A \to B\) 为 Noether 局部环之间的局部同态，M 为有限生成 A-模，N 为有限生成 B-模；我们假设 A-模 N 平坦。

\begin{enumerate}
\def\labelenumi{\alph{enumi})}
\item
  设 \((x_{1},\ldots,x_{r})\) 为 \(m_{A}\) 中元素组成的、对 A-模 M 正则的序列，\((y_{1},\ldots,y_{s})\) 为 \(m_{B}\) 中元素组成的、对 B-模 \(\kappa_{A}\otimes_{A}N\) 正则的序列；则序列 \((y_{1},\ldots,y_{s},\rho(x_{1}),\ldots,\rho(x_{r}))\) ——由 \(m_{B}\) 中元素组成——对 B-模 \(M\otimes_{A}N\) 正则。
\item
  我们有等式
\end{enumerate}

\[
\mathrm{prof} _ {\mathrm{B}} (\mathrm{M} \otimes_ {\mathrm{A}} \mathrm{N}) = \mathrm{prof} _ {\mathrm{A}} (\mathrm{M}) + \mathrm{prof} _ {\mathrm{B}} (\kappa_ {\mathrm{A}} \otimes_ {\mathrm{A}} \mathrm{N}).
\]

设 \(\mathfrak{x}\) 为由序列 \(\mathbf{x}\) 生成的 A 的理想，\(\mathfrak{y}\) 为由 \(y\) 生成的 B 的理想。由命题 10，序列 \(y\) 是 \(M \otimes_A N\) -正则的，且 \(N/\mathfrak{y}N\) 在 \(A\) 上平坦，于是序列 \(\boldsymbol{\rho}(\mathbf{x}) = (\boldsymbol{\rho}(x_1), \ldots, \boldsymbol{\rho}(x_r))\) 对 \(M \otimes_A (N/\mathfrak{y}N) = (M \otimes_A N)/\mathfrak{y}(M \otimes_A N)\) 正则。这就证明了 a)。

为证明 b)，我们可设 M 与 N 非零。由 Nakayama 引理，\(\kappa_{\mathrm{A}} \otimes_{\mathrm{A}} \mathrm{N}\) 也非零，故 \(\operatorname{prof}_{\mathrm{A}}(\mathrm{M})\) 与 \(\operatorname{prof}_{\mathrm{B}}(\kappa_{\mathrm{A}} \otimes_{\mathrm{A}} \mathrm{N})\) 有限（第 4 节, 定理 2 的推论 2）。取极大的正则序列 x 与 y；于是有 \(r = \operatorname{prof}_{\mathrm{A}}(\mathrm{M})\) ， \(s = \operatorname{prof}_{\mathrm{B}}(\kappa_{\mathrm{A}} \otimes_{\mathrm{A}} \mathrm{N})\) ，且存在单 A-线性映射 \(u: \kappa_{\mathrm{A}} \to \mathrm{M} / x \mathrm{M}\) 与单 B-线性映射 \(v: \kappa_{\mathrm{B}} \to \kappa_{\mathrm{A}} \otimes_{\mathrm{A}} (\mathrm{N} / \mathfrak{y} \mathrm{N})\) 。由于 N/ \(\mathfrak{y}\) N 在 A 上平坦，B-线性映射 \((u \otimes 1_{\mathrm{N} / \mathfrak{y} \mathrm{N}}) \circ v\) ——它把 \(\kappa_{\mathrm{B}}\) 映入 (M/xM) \(\otimes_{\mathrm{A}} (\mathrm{N} / \mathfrak{y} \mathrm{N}) = (\mathrm{M} \otimes_{\mathrm{A}} \mathrm{N}) / (\rho(x) + \mathfrak{y})(\mathrm{M} \otimes_{\mathrm{A}} \mathrm{N})\) ——是单射的。这蕴含等式 \(\operatorname{prof}_{\mathrm{B}}(\mathrm{M} \otimes_{\mathrm{A}} \mathrm{N}) = r + s\) （同上引文）。

注. -- 回忆在此前诸假设下，我们有

\[
\dim_ {B} (M \otimes_ {A} N) = \dim_ {A} (M) + \dim_ {B} (\kappa_ {A} \otimes_ {A} N)
\]

(VIII, § 3, 第 4 节, 命题 7)。

推论.--- 设 \(\rho: A \to B\) 为 Noether 局部环之间的使 B 成为平坦 A-模的局部同态。我们有

\[
\begin{array}{l} \operatorname{prof} (\mathrm{B}) = \operatorname{prof} (\mathrm{A}) + \operatorname{prof} \left(\kappa_ {\Lambda} \otimes_ {\Lambda} \mathrm{B}\right), \\ \dim (\mathrm{B}) = \dim (\mathrm{A}) + \dim \left(\kappa_ {\mathrm{A}} \otimes_ {\mathrm{A}} \mathrm{B}\right). \end{array}
\]

事实上，由命题 4 的推论（相应地，由 VIII, § 1, 第 4 节），B-模 \(\kappa_{\mathrm{A}}\otimes_{\mathrm{A}}\mathrm{B}\) 的深度（相应地，维数）等于环 \(\kappa_{\mathrm{A}}\otimes_{\mathrm{A}}\mathrm{B}\) 的深度（相应地，维数）。

\subsection{深度的界}

命题 12.--- 设 A 为 Noether 局部环，M 为非零的有限生成 A-模，J 为异于 A 的 A 的理想。我们有不等式链

\[
\begin{array}{r l} \operatorname{prof} _ {A} (J; M) & \leqslant \operatorname{codim} (\operatorname{Supp} (M) \cap V (J), \operatorname{Supp} (M)) \leqslant \dim (M) - \dim (M / J M) \\ & \leqslant [ J / \mathfrak {m} _ {A} J: \kappa_ {A} ]. \end{array}
\]

对每个 \(\mathfrak{p} \in \operatorname{Supp}(\mathrm{M})\cap \mathrm{V}(\mathrm{J})\) ，\(\operatorname{prof}_{\mathrm{A}}(\mathrm{J};\mathrm{M})\) 小于 \(\dim_{\mathrm{A}_{\mathfrak{p}}}(\mathrm{M}_{\mathfrak{p}})\) （第 5 节, 命题 8 及第 4 节, 定理 2 的推论 2），即 (VIII, § 1, 第 4 节, 命题 9)

即 \(\operatorname{codim}(V(\mathfrak{p}),\operatorname{Supp}(\mathrm{M}))\) 。当 \(\mathfrak{p}\) 遍历 \(\operatorname{Supp}(\mathrm{M}) \cap \mathrm{V}(\mathrm{J})\) 时，\(V(\mathfrak{p})\) 遍历 \(\operatorname{Supp}(\mathrm{M}) \cap \mathrm{V}(\mathrm{J})\) 的不可约闭子集，由此得第一个不等式。第二个不等式由 VIII, § 1, 第 2 节的命题 3 得出。此外，可以找到 J 的一个基数为 \([J/\mathfrak{m}_{\mathrm{A}}J:\kappa_{\mathrm{A}}]\) 的生成组 (II, § 3, 第 2 节, 命题 4 的推论 2)；第三个不等式于是由 VIII, § 3, 第 2 节, 公式 (8) 得出。

注 1.-- 考察命题 12 中的不等式链。

\begin{enumerate}
\def\labelenumi{\alph{enumi})}
\item
  要使 \(\mathrm{prof}_{\mathrm{A}}(\mathrm{J};\mathrm{M}) = [\mathrm{J} / \mathfrak{m}_{\mathrm{A}}\mathrm{J}:\kappa_{\mathrm{A}}]\) 成立，必须且只需 J 可由一个 M-正则序列生成（第 3 节, 定理 1 的推论 2 及 A, X, 第 160 页, 推论 1）。
\item
  等式 \(\dim(\mathrm{M}) - \dim(\mathrm{M}/\mathrm{JM}) = [\mathrm{J}/\mathfrak{m}_{\Lambda}\mathrm{J} : \kappa_{\Lambda}]\) 意味着 J 可由一个对 M 正则的序列生成 (VIII, § 3, 第 2 节)。
\item
  若 M 是 Macaulay 的，则有 \(\operatorname{prof}_{\mathrm{A}}(\mathbf{J};\mathbf{M}) = \dim (\mathbf{M}) - \dim (\mathbf{M} / \mathbf{JM})\) (§ 2, 第 2 节, 命题 2 的推论)。
\end{enumerate}

引理 3.--- 设 A 为 Noether 环，\(\mathfrak{p} \subset \mathfrak{p}_1 \subset \ldots \subset \mathfrak{p}_{r-1} \subset \mathfrak{q}\) 为 A 的素理想组成的长度为 \(r\) 的饱和链，M 为有限型 \(\Lambda\) -模，\(n\) 为整数。若 A-模 \(\mathrm{Ext}_{\mathrm{A}_{\mathfrak{p}}}^{n}(\kappa(\mathfrak{p}), \mathrm{M}_{\mathfrak{p}})\) 非零，则 \(\mathrm{Ext}_{\mathrm{A}_{\mathfrak{q}}}^{n+r}(\kappa(\mathfrak{q}), \mathrm{M}_{\mathfrak{q}})\) 亦然。

显然只需处理 \(r = 1\) 的情形；然后把 A、M、p 与 q 分别用 \(A_q\) 、 \(M_q\) 、 \(pA_q\) 与 \(qA_q\) 代换，便归结为处理 A 为局部环且 \(q = m_A\) 的情形。设 \(x\) 为 \(m_A - p\) 的元素。\(A_p\) -模 \(\text{Ext}_A^n(A/p, M) \otimes_A A_p\) 同构于 \(\text{Ext}_{A_p}^n(\kappa(p), M_p)\) (A, X, 第 111 页, 命题 10 b))，故按假设非零；更有 \(\text{Ext}_A^n(A/p, M)\) 不为零。正合序列

\[
0 \rightarrow \mathrm{A} / \mathfrak {p} \stackrel {{x _ {\mathrm{A} / \mathfrak {p}}}} {{\longrightarrow}} \mathrm{A} / \mathfrak {p} \longrightarrow \mathrm{A} / (\mathfrak {p} + x \mathrm{A}) \rightarrow 0
\]

给出扩张模的一个正合序列

\[
\operatorname{Ext} _ {\mathrm{A}} ^ {n} (\mathrm{A} / \mathfrak {p}, \mathrm{M}) \xrightarrow {u} \operatorname{Ext} _ {\mathrm{A}} ^ {n} (\mathrm{A} / \mathfrak {p}, \mathrm{M}) \longrightarrow \operatorname{Ext} _ {\mathrm{A}} ^ {n + 1} (\mathrm{A} / (\mathfrak {p} + x \mathrm{A}), \mathrm{M}),
\]

其中 \(u\) 是乘以 \(x\) 的映射。由 Nakayama 引理，它不是满射，故 A-模 \(\text{Ext}_A^{n+1}(\text{A}/(\mathfrak{p} + x\text{A}), \text{M})\) 不为零。而 A 的包含 \(\mathfrak{p} + x\text{A}\) 的素理想只有 \(\mathfrak{m}_\text{A}\) ；因此 A-模 \(\text{A}/(\mathfrak{p} + x\text{A})\) 有有限长度 (VIII, § 3, 第 2 节, 引理 2)。若 \(\text{Ext}_\text{A}^{n+1}(\kappa_\text{A}, \text{M})\) 为零，则由对 N 的长度作归纳将推出，对每个有限长度的 A-模 N，\(\text{Ext}_\text{A}^{n+1}(\text{N}, \text{M})\) 为零。这一矛盾完成了证明。

命题 13.--- 设 A 为 Noether 环，M 为有限生成 A-模，p 与 q 为 Supp(M) 中满足 p ⊂ q 的两个素理想。则

\[
\operatorname{prof} _ {A _ {\mathfrak {q}}} (M _ {\mathfrak {q}}) \leqslant \operatorname{prof} _ {A _ {\mathfrak {p}}} (M _ {\mathfrak {p}}) + \dim (A _ {\mathfrak {q}} / \mathfrak {p} A _ {\mathfrak {q}}).
\]

更精确地说，对素理想的每个饱和链 \(\mathfrak{p} \subset \mathfrak{p}_1 \subset \ldots \subset \mathfrak{p}_{r-1} \subset \mathfrak{q}\) ，有 \(\operatorname{prof}_{\mathrm{A}_{\mathfrak{q}}}(\mathrm{M}_{\mathfrak{q}}) \leqslant \operatorname{prof}_{\mathrm{A}_{\mathfrak{p}}}(\mathrm{M}_{\mathfrak{p}}) + r\) 。

令 \(p = \text{prof}_{\text{A}_p}(\text{M}_p)\) ，我们来证明第二个不等式。若 \(p = +\infty\) 它是显然的；否则，有 \(\text{Ext}_{\text{A}_p}^p(\kappa(\mathfrak{p}), \text{M}_p) \neq 0\) ，从而由引理 3 得 \(\operatorname{Ext}_{\Lambda_{\mathfrak{q}}}^{p+r}(\kappa(\mathfrak{q}),\mathrm{M}_{\mathfrak{q}})\neq0\) ，这蕴含 \(\operatorname{prof}_{\mathrm{A}_{\mathfrak{q}}}(M_{\mathfrak{q}})\leqslant p+r\) 。由于 \(\dim(A_{\mathfrak{q}}/\mathfrak{p}A_{\mathfrak{q}})\) 是以给定端点 \(\mathfrak{p}\) 与 \(\mathfrak{q}\) 的素理想饱和链的长度的上确界，第一个断言是第二个断言的推论。

推论 1.--- 我们有不等式

\[
\dim (M _ {\mathfrak {q}}) - \operatorname{prof} _ {\Lambda_ {\mathfrak {q}}} (M _ {\mathfrak {q}}) \geqslant \dim (M _ {\mathfrak {p}}) - \operatorname{prof} _ {A _ {\mathfrak {p}}} (M _ {\mathfrak {p}}) \geqslant 0.
\]

这由命题 13 及不等式 \(\dim(M_{\mathfrak{q}}) \geqslant \dim(M_{\mathfrak{p}}) + \dim(A_{\mathfrak{q}} / \mathfrak{p}A_{\mathfrak{q}})\) (VIII, § 1, 第 4 节, 命题 9, a) 及第 3 节, 命题 7, b)) 得出。

推论 2.--- 设 A 为 Noether 局部环，M 为有限生成 A-模。则有不等式

\[
\operatorname{prof} _ {\mathrm{A}} (\mathrm{M}) \leqslant \inf _ {\mathfrak {p} \in \operatorname{Ass} (\mathrm{M})} \dim (\mathrm{A} / \mathfrak {p}).
\]

设 \(\mathfrak{p}\) 为与 M 相伴的素理想；有 \(\mathrm{prof}_{\mathrm{A}_{\mathfrak{p}}}(\mathrm{M}_{\mathfrak{p}})=0\) （第 1 节, 注 2）。把命题 13 应用于理想 \(\mathfrak{p}\subset\mathfrak{m}_{\mathrm{A}}\) 便给出不等式 \(\mathrm{prof}_{\mathrm{A}}(\mathrm{M})\leqslant\dim(\mathrm{A}/\mathfrak{p})\) ，从而得本推论。

注 2.--- 我们有 \(\sup_{\mathfrak{p}\in\operatorname{Ass}(\mathbb{M})}\dim(\mathrm{A}/\mathfrak{p})=\dim(\mathrm{M})\) (VIII, § 1, 第 4 节, 注 2)，并且重新得到不等式 \(\operatorname{prof}(\mathrm{M})\leqslant\dim(\mathrm{M})\) （对 \(\mathrm{M}\neq0\) ）（第 4 节, 定理 2 的推论 2）。

\subsection{局部整 Noether 环；正规 Noether 环}

设 A 为 Noether 环。用 \((\mathrm{Y}_{i})_{i\in\mathrm{I}}\) 表示 Spec(A) 的连通分支组成的有限族 (II, § 4, 第 2 节, 命题 10 及第 3 节, 命题 11 的推论 7)。根据 II, § 4, 第 3 节, 命题 15，对每个 i 存在 A 的唯一幂等元 \(e_{i}\) 使 \(\mathrm{Y}_{i}=\mathrm{V}(e_{i})\) ，且从 A 到诸环 A/Ae \(_{i}\) 之积的典范映射是双射。A 的商环 A/Ae \(_{i}\) 称为 A 的典范分量。令 \(f_{i}=1-e_{i}\) 。有 \(\sum_{i\in I}f_{i}=1\) ，且 \((f_{i})_{i\in I}\) 是 A 的非零幂等元组成的正交族（同上引文）。由此得出，\(f_{i}\) 在 A/Ae \(_{j}\) 中的像当 j=i 时等于 1，当 j≠i 时等于 0；于是由环同态 A→ \(\prod_{j}A/Ae_{j}\) 得出典范同构 A \(_{f_{i}}\rightarrow A/Ae_{i}\) 。

根据同上引文, 命题 14 的推论 2，下列条件等价： Spec(A) 的连通分支是不可约的；

\begin{enumerate}
\def\labelenumi{(\roman{enumi})}
\setcounter{enumi}{1}
\item
  A 的每个素理想（相应地，极大理想）只属于 Spec(A) 的一个不可约分支；
\item
  A 的每个素理想（相应地，极大理想）只包含一个极小素理想；
\item
  对 A 的每个素理想（相应地，极大理想）p，拓扑空间 \(\operatorname{Spec}(A_{\mathfrak{p}})\) 是不可约的；
\item
  对 A 的每个典范分量 C，拓扑空间 \(\operatorname{Spec}(\mathbb{C})\) 是不可约的。
\end{enumerate}

现在注意，若 A 是约化的，则所有环 \(A_{p}\) 都是约化的 (II, § 2, 第 6 节, 命题 17)；反之，若对 A 的每个极大理想 m，环 \(A_{m}\) 都约化，则 A 约化 (II, § 3, 第 3 节, 定理 1 的推论 2)。于是应用 II, § 4, 第 3 节, 命题 14 的推论 1，由前面所述推出下列条件的等价性：

\begin{enumerate}
\def\labelenumi{(\roman{enumi})}
\item
  A 约化且 Spec(A) 的连通分支不可约；
\item
  对 A 的每个素理想（相应地，极大理想）p，环 A \(_{p}\) 是整环；
\item
  A 的典范分量都是整环。
\end{enumerate}

若一个 Noether 环满足上述等价条件 至 ，就称它是局部整的。

设环 A 局部整；设 u 为 A 到整环的（有限）积 \(\prod_{j\in J}A_{j}\) 上的同构。存在双射 \(\sigma:J\to I\) ，使与 u 相伴的从 \(\operatorname{Spec}(\prod_{j\in J}A_{j})\) 到 \(\operatorname{Spec}(A)\) 的映射给出 \(\operatorname{Spec}(A_{j})\) 到连通分支 \(Y_{\sigma(j)}\) —— \(\operatorname{Spec}(A)\) 的连通分支——上的同胚；于是由 u 得出典范分量 \(A/Ae_{\sigma(j)}\) 到 \(A_{j}\) 上的同构。

命题 14. - 设 A 为 Noether 环。下列条件等价：

\begin{enumerate}
\def\labelenumi{(\roman{enumi})}
\item
  A 约化且在其全分式环中整闭；
\item
  A 同构于一个有限族整闭环的积。
\item
  A 的典范分量是整闭的；
\item
  对 A 的每个素理想（相应地，极大理想）p，环 \(A_{p}\) 整闭。
\end{enumerate}

 与 的等价性由 V, § 1, 第 2 节, 命题 9 的推论 2 得出， 与 的等价性由命题前的注记得出。设 p 为 A 的素理想；存在 A 的唯一典范分量 A' 使 p 属于 Spec(A) 的闭子集 Spec(A')，且我们有典范同构 \(A_{p} \simeq A'_{p}\)。于是 与 的等价性由同上引文, 第 5 节, 命题 16 的推论 1 与推论 3 得出。

若环 A 是 Noether 的且满足命题 14 的等价条件 至 ，则称 A 为正规环。Noether 环整闭当且仅当它是整环且正规。正规局部环是整闭的。

\subsection{深度与连通性}

引理 4.--- 设 A 为 Noether 环，F 为 Spec(A) 的闭子集，U 为其补开集，u : M → N 为有限生成 A-模之间的同态。

\begin{enumerate}
\def\labelenumi{\alph{enumi})}
\item
  设 \(u_{\mathfrak{p}}: \mathrm{M}_{\mathfrak{p}} \to \mathrm{N}_{\mathfrak{p}}\) 对所有 \(\mathfrak{p} \in \mathrm{U}\) 单射，并设 \(\operatorname{prof}_{\mathrm{F}}(\mathrm{M}) \geqslant 1\) 。则 \(u\) 单射。
\item
  设 \(u_{\mathfrak{p}}: \mathrm{M}_{\mathfrak{p}} \to \mathrm{N}_{\mathfrak{p}}\) 对所有 \(\mathfrak{p} \in \mathrm{U}\) 双射，并设 \(\operatorname{prof}_{\mathrm{F}}(\mathrm{M}) \geqslant 2\) 与 \(\operatorname{prof}_{\mathrm{F}}(\mathrm{N}) \geqslant 1\) 。则 u 双射。
\item
  诸假设蕴含 Supp(Ker \(u\) ) \(\subset\) F，于是 \(\operatorname{Hom}_{\Lambda}(\operatorname{Ker} u, \mathbf{M}) = 0\) （ \(n^{\circ}\) 5, 注 1）；因此有 \(\operatorname{Ker} u = 0\) 。
\item
  我们已知 \(u\) 单射，且有 \(\operatorname{Supp}(\operatorname{Coker} u) \subset \mathbb{F}\) 。由同上引文，有 \(\operatorname{Hom}_{\mathbf{A}}(\operatorname{Coker} u, \mathbf{N}) = 0\) 与 \(\operatorname{Ext}_{\mathbf{A}}^{1}(\operatorname{Coker} u, \mathbf{M}) = 0\) 。由扩张模的正合序列
\end{enumerate}

\[
\operatorname{Hom} _ {\mathrm{A}} (\operatorname{Coker} u, \mathrm{N}) \rightarrow \operatorname{Hom} _ {\mathrm{A}} (\operatorname{Coker} u, \operatorname{Coker} u) \rightarrow \operatorname{Ext} _ {\mathrm{A}} ^ {1} (\operatorname{Coker} u, \mathrm{M})
\]

我们推出 \(\operatorname{Hom}_{\mathrm{A}}(\operatorname{Coker} u, \operatorname{Coker} u) = 0\) ，这蕴含 \(\operatorname{Coker} u = 0\) 。

注 1.—设 A 为 Noether 环，F 为 \(\operatorname{Spec}(A)\) 的闭子集，U 为其补开集。要使 \(\operatorname{prof}_{\mathrm{F}}(A) \geqslant 1\) ，必须且只需 \(\operatorname{Ass}(A) \subset U\) （注 2, 第 5 节）。当此条件满足时，\(\operatorname{Spec}(A)\) 的每个不可约分支都与 U 相交，故 U 在 \(\operatorname{Spec}(A)\) 中稠密。

定理 3 (Hartshorne).--- 设 A 为 Noether 环，F 为 Spec(A) 的闭子集，U 为其补开集。设 prof \(_{F}\) (A) ≥ 2。则对 Spec(A) 的每个连通分支 Y，集合 Y ∩ U 连通且在 Y 中稠密。

首先设 \(\operatorname{Spec}(A)\) 连通。由注 1，\(U\) 在 \(\operatorname{Spec}(A)\) 中稠密，剩下要证明它连通。反设两个不相交的非空开子集 \(U_0\) 与 \(U_1\) —— \(\operatorname{Spec}(A)\) 的开子集——之并为 \(U\)。由注 1，\(\operatorname{Ass}(A)\) 包含于 \(U\)，故它是 \(\operatorname{Ass}(A) \cap U_0\) 与 \(\operatorname{Ass}(A) \cap U_1\) 的不交并。由 IV, § 1, 第 1 节, 命题 4，存在 A 的理想 \(J_0\) 与 \(J_1\) 使 \(\operatorname{Ass}(J_i) = \operatorname{Ass}(A) \cap U_i\) ， \(\operatorname{Ass}(A/J_i) = \operatorname{Ass}(A) \cap U_{1-i}\quad (i=0,1)\) 。\(U_i\) 在 \(\operatorname{Spec}(A)\) 中的补集包含 \(\operatorname{Ass}(A/J_i)\) 与 \(\operatorname{Ass}(J_{1-i})\)；由于它是闭集，它包含 \(\operatorname{Supp}(A/J_i)\) 与 \(\operatorname{Supp}(J_{1-i})\)。于是对 \(\mathfrak{p} \in U_i\) ，有 \((J_i)_{\mathfrak{p}} = A_{\mathfrak{p}}\) 与 \((J_{1-i})_{\mathfrak{p}} = 0\) ；特别地这蕴含 \(J_0\) 与 \(J_1\) 异于 A。设 B 为 A-模 \(A/J_0 \times A/J_1\) ，设 \(u: A \to B\) 为典范同态。由前所述，同态 \(u_{\mathfrak{p}}\) 对每个 \(\mathfrak{p} \in U\) 双射；此外有 \(\operatorname{Ass}(B) \subset U_0 \cup U_1 = U\) ，故由注 1 有 \(\operatorname{prof}_F(B) \geqslant 1\) 。于是引理 4 蕴含 \(u\) 双射，这与 \(\operatorname{Spec}(A)\) 的连通性矛盾。

现在处理一般情形。设 J 为 A 的理想使 \(F = V(J)\)，并设 Y 为 \(\operatorname{Spec}(A)\) 的一个连通分支。由 II, § 4, 第 3 节, 命题 15，存在 A 的幂等元 \(f\) 使 Y 等同于 \(\operatorname{Spec}(A_f)\) （ \(\operatorname{Spec}(A)\) 的一个部分）。于是 \(Y \cap F\) 等同于 \(V(J_f)\)；由第 3 节的命题 6 a)，有 \(\operatorname{prof}_{A_f}(J_f, A_f) \geqslant \operatorname{prof}_A(J; A) \geqslant 2\) 。由证明的第一部分得出，\(Y \cap U = Y - (Y \cap F)\) 连通且在 Y 中稠密。

推论 1.--- 把 U 的每个连通分支映为其在 Spec(A) 中的闭包的映射，是从 U 的连通分支组成的集合到 Spec(A) 的连通分支组成的集合上的一个双射。

推论 2.--- 对每个深度 \(\geqslant\) 2 的 Noether 局部环 B，拓扑空间 \(\operatorname{Spec}(\mathbf{B}) = \{\mathfrak{m}_{\mathbf{B}}\}\) 连通。

推论 3.--- 在定理 3 的假设下，设 \(\text{Spec}(A_{\mathfrak{p}})\) 不可约（相应地，\(A_{p}\) 为整环）对一切 \(p \in U\) 成立；则 \(\text{Spec}(A_{\mathfrak{p}})\) 不可约（相应地，\(A_{p}\) 为整环）对一切 \(p \in \text{Spec}(A)\) 成立。

设 \((\mathrm{Y}_{i})_{i\in\mathrm{I}}\) 为 \(\operatorname{Spec}(A)\) 的不可约分支组成的（有限）族。设 \(\mathfrak{p}\in U\) ；由于 \(\operatorname{Spec}(A_{\mathfrak{p}})\) 不可约，\(\mathfrak{p}\) 只包含 A 的一个极小素理想，因而只属于诸 \(Y_i\) 之一 (II, § 4, 第 3 节, 命题 14 的推论 2)。子集 \(Y_i\cap U\) 是 U 的闭子集，两两不交，且因 U 在 \(\operatorname{Spec}(A)\) 中稠密而非空，并由 II, § 4, 第 1 节, 命题 7 不可约；因而它们是 U 的连通分支。它们的闭包 \(Y_i\) 是 \(\operatorname{Spec}(A)\) 的连通分支（推论 1）。这就证明了 \(\operatorname{Spec}(A)\) 的连通分支不可约，从而 \(\operatorname{Spec}(A_{\mathfrak{p}})\) 对每个 \(\mathfrak{p}\) 不可约（第 8 节）。

设 \(A_{\mathfrak{q}}\) 对所有 \(\mathfrak{q} \in U\) 是整的。设 \(\mathfrak{p} \in \operatorname{Spec}(A)\) 。由于 \(\operatorname{Spec}(A_{\mathfrak{p}})\) 不可约，\(A_{\mathfrak{p}}\) 的诣零根是 \(A_{\mathfrak{p}}\) 的唯一极小素理想；因此它属于 \(\operatorname{Ass}(A_{\mathfrak{p}})\) (IV, § 1, 第 3 节, 命题 7 的推论 1)，从而等于 \(\mathfrak{q}A_{\mathfrak{p}}\) ，其中 \(\mathfrak{q}\) 是与 A 相伴的素理想 (IV, § 1, 第 2 节, 命题 5 的推论)。有 \(\mathfrak{q} \in U\) （注 1）且 \(\mathfrak{q}A_{\mathfrak{q}} \in \operatorname{Ass}(A_{\mathfrak{q}})\) （同上引文）；由于 \(A_{\mathfrak{q}}\) 是整的，\(\mathfrak{q}\) 为零，因此 \(A_{\mathfrak{p}}\) 是整环。

推论 4.--- 设 \(A\) 为谱连通的 Noether 环。设存在整数 \(d \geqslant 1\) ，使 \(\operatorname{prof}(A_{\mathfrak{p}}) \geqslant 2\) 对 A 的每个素理想 \(\mathfrak{p}\) （高度 \textgreater{} d 者）成立。

\begin{enumerate}
\def\labelenumi{\alph{enumi})}
\item
  对 \(\operatorname{Spec}(\mathbf{A})\) 的每个余维数 \(>d\) 的闭子集 Z，空间 \(\operatorname{Spec}(\mathbf{A}) - \mathbf{Z}\) 连通。
\item
  设 Y 与 Y' 为 Spec(A) 的不可约分支。存在 Spec(A) 的不可约分支的序列 \(X_{1}, X_{2}, \ldots, X_{n}\) ，使得有 \(X_{1} = Y\) ， \(X_{n} = Y'\) ，且对 \(i = 1, \ldots, n - 1\)
\end{enumerate}

\[
\operatorname{codim} \left(\mathrm{X} _ {i} \cap \mathrm{X} _ {i + 1}, \operatorname{Spec} (\mathrm{A})\right) \leqslant d.
\]

设 \(Z \subset \operatorname{Spec}(A)\) 为余维数 \(>d\) 的闭子集。对每个 \(\mathfrak{p} \in Z\) ，有 \(\dim(A_{\mathfrak{p}}) > d\) ，从而 \(\operatorname{prof}(A_{\mathfrak{p}}) \geqslant 2\) ，这蕴含 \(\operatorname{prof}_{Z}(A)\) 为 \(\geqslant 2\)（第 5 节, 命题 8），从而 \(\operatorname{Spec}(A) - Z\) 连通（定理 3）。

设 Z 为诸集合 \(X^{\prime} \cap X^{\prime\prime}\) 之并，其中 \((X^{\prime}, X^{\prime\prime})\) 遍历 Spec(A) 的不可约分支的（有限）对组成的集合中满足 \(\text{codim}(X^{\prime} \cap X^{\prime\prime}, \text{Spec}(A)) > d\) 者。由 a)，集合 \(\text{Spec}(A) - Z\) 连通。\(\text{Spec}(A)\) 的所有不可约分支都与 \(\text{Spec}(A) - Z\) 相交；它们在 \(\text{Spec}(A) - Z\) 上的迹是 \(\text{Spec}(A) - Z\) 的不可约分支 (II, § 4, 第 1 节, 命题 7)。由于 \(\text{Spec}(A) - Z\) 连通，存在序列 \(X_{1}, \ldots, X_{n}\) ，由 \(\text{Spec}(A)\) 的不可约分支组成，使得有 \(X_{1} - Z = Y - Z\) ， \(X_{n} - Z = Y' - Z\) ，且 \((X_{i} - Z) \cap (X_{i+1} - Z) \neq \emptyset\) 对 \(1 \leqslant i \leqslant n - 1\) 成立；换言之，有 \(X_{1} = Y\) ， \(X_{n} = Y'\) 且 \(\text{codim}(X_{i} \cap X_{i+1}) \leqslant d\) 。

注 2. 当 A 是 Macaulay 环时，可对 \(d=1\) 应用该推论（ \(\S\) 2, 第 5 节）。

例（由 4 维仿射空间的四个坐标平面构成的完全交）.--- 设 \(k\) 为域。用 S 表示多项式环 \(k[\mathrm{T}_{1},\mathrm{T}_{2},\mathrm{T}_{3},\mathrm{T}_{4}]\)。回忆 (VIII, § 2, 第 4 节, 定理 3)，S 的素理想的每个极大链长度为 4。设 \(\mathfrak{m}\) 为 S 的由诸 \(\mathrm{T}_{i}\) 生成的理想，\(\mathfrak{a}\) 为由 \(\mathrm{T}_{1}\mathrm{T}_{2}\) 与 \(\mathrm{T}_{3}\mathrm{T}_{4}\) 生成的理想，且对 \(1 \leqslant i < j \leqslant 4\) ，设 \(\mathfrak{p}_{ij}\) 为由 \(\mathrm{T}_{i}\) 与 \(\mathrm{T}_{j}\) 生成的理想。诸理想 \(\mathfrak{p}_{ij}\) 是高度 2 的素理想，它们的和是极大理想 \(\mathfrak{m}\) ，且有 \(\mathfrak{a} = \mathfrak{p}_{13} \cap \mathfrak{p}_{14} \cap \mathfrak{p}_{23} \cap \mathfrak{p}_{24}\)。

\begin{enumerate}
\def\labelenumi{\alph{enumi})}
\item
  环 A = S/a 是约化的；Spec(A) 的不可约分支是诸集合 \(X_{ij} = V(\mathfrak{p}_{ij}/\mathfrak{a})\) （i = 1, 2, j = 3, 4），它们维数为 2 且都包含点 m/a。特别地，Spec(A) 连通且维数为 2。两个不同分支 \(X_{ij}\) 与 \(X_{kl}\) 的交退化为 \(\{m/a\}\) 当 \(\{i,j\} \cap \{k,l\} = \varnothing\) 时，否则维数为 1。由此得出，推论 4 的结论对 d = 1 成立（我们后面将看到 A 是 Macaulay 环，因而推论 4 的假设本身对 d = 1 成立）。
\item
  设 b 为 S 的由 \(T_{1}T_{2}\) 、 \(T_{1}T_{3}\) 、 \(T_{2}T_{4}\) 、 \(T_{3}T_{4}\) 生成的理想，B 为环 S/b。有 \(b = p_{14}p_{23} = p_{14} \cap p_{23}\) 。环 B 是约化的。它的谱等同于 Spec(A) 的闭子集 \(X_{14} \cup X_{23}\) ；它有两个不可约分支（维数为 2），其交退化为一点。B 沿这一点的深度因 B 约化而严格为正，且由定理 3 知小于 2，故等于 1。于是 B 不是 Macaulay 环。
\end{enumerate}

\subsection{深度与正规性}

设 A 为 Noether 环，p 为 A 的素理想。有 \(\text{prof}(A_{\mathfrak{p}}) \leqslant \text{ht}(\mathfrak{p})\) （第 4 节, 定理 2 的推论 2）。此外若 A 约化，则局部环 \(A_{p}\) 约化 (II, § 2, 第 6 节, 命题 17)，由此蕴含：

\begin{enumerate}
\def\labelenumi{\alph{enumi})}
\item
  若 ht(p) = 0，则 A \(_{p}\) 是域；
\item
  若 \(\operatorname{ht}(\mathfrak{p}) \geqslant 1\) ，则 \(\operatorname{prof}(A_{\mathfrak{p}}) \geqslant 1\) （ \(n^{\circ}\) 1, 注 3）。
\end{enumerate}

反之：

命题 15.--- 设 A 为满足下列两个条件的 Noether 环：

\begin{enumerate}
\def\labelenumi{(\roman{enumi})}
\item
  对 A 的每个极小素理想 \(\mathfrak{p}\) ，环 \(A_{\mathfrak{p}}\) 约化；
\item
  对 A 的每个素理想 \(\mathfrak{p}\) （高度 \(\geqslant 1\) 者），有 \(\operatorname{prof}(A_{\mathfrak{p}}) \geqslant 1\) 。则 A 约化。
\end{enumerate}

设 \(\mathfrak{n}\) 为 A 的诣零根。对 A 的每个极小素理想 \(\mathfrak{p}\) ，由 (i) 有 \(\mathfrak{n}_{\mathfrak{p}} = 0\) ，也就是 \(\mathfrak{p} \notin \operatorname{Supp}(\mathfrak{n})\) ，更有 \(\mathfrak{p} \notin \operatorname{Ass}_{\Lambda}(\mathfrak{n})\) 。对每个 \(\mathfrak{p} \in \operatorname{Spec}(\mathbf{A})\) （高度 \(\geqslant 1\) 者），由 有 \(\mathfrak{pA}_{\mathfrak{p}} \notin \operatorname{Ass}_{\mathrm{A}_{\mathfrak{p}}}(\mathrm{A}_{\mathfrak{p}})\) ，更有 \(\mathfrak{pA}_{\mathfrak{p}} \notin \operatorname{Ass}_{\mathrm{A}_{\mathfrak{p}}}(\mathfrak{n}_{\mathfrak{p}})\) ，从而 \(\mathfrak{p} \notin \operatorname{Ass}_{\mathrm{A}}(\mathfrak{n})\) (IV, § 1, 第 2 节, 命题 5)。于是集合 \(\operatorname{Ass}_{\mathrm{A}}(\mathfrak{n})\) 为空，这蕴含 \(\mathfrak{n}\) 为零。

命题 16.--- 设 A 为整闭的 Noether 环，J 为 A 的高度 ≥ 2 的理想，M 为自反的有限生成 A-模。则 prof \(_{A}\) (J; M) ≥ 2。

选取 A 的分式域上的一个有限维向量空间 V，以及 V 中同构于 M 的一个格 N (VII, § 4, 第 2 节, 注 1)。V/N 的相伴素理想高度为 1（同上引文, 定理 2），故不包含 J；由第 1 节的注 2，这蕴含 \(\operatorname{prof}_{\mathrm{A}}(\mathrm{J};\mathrm{V}/\mathrm{N}) \geqslant 1\) 。另一方面，A-模 V 是可除的且无挠，故为内射的 (A, X, 第 17 页, 命题 10 的推论 2)，这蕴含 \(\operatorname{prof}_{\mathrm{A}}(\mathrm{J};\mathrm{V}) = +\infty\) 。于是不等式 \(\operatorname{prof}_{\mathrm{A}}(\mathrm{J};\mathrm{N}) \geqslant 2\) 由第 1 节的命题 1 得出。

推论.--- 整闭且维数 ≥ 2 的 Noether 局部环深度 ≥ 2。

设 A 为正规环（第 8 节），p 为 A 的素理想。则：

\begin{enumerate}
\def\labelenumi{\alph{enumi})}
\item
  若 ht(p) = 0，则 A \(_{p}\) 是域；
\item
  若 \(\operatorname{ht}(\mathfrak{p}) = 1\) ，则 \(A_{p}\) 是离散赋值环 (VII, § 1, 第 7 节, 命题 11)；
\item
  若 \(\operatorname{ht}(\mathfrak{p}) \geqslant 2\)，则 \(\operatorname{prof}(A_{\mathfrak{p}}) \geqslant 2\) （命题 16 的推论）。
\end{enumerate}

反之：

定理 4 (Serre).--- 设 A 为满足下列条件的 Noether 环：

\begin{enumerate}
\def\labelenumi{(\roman{enumi})}
\item
  对 A 的每个极小素理想 p，环 A \(_{p}\) 约化；
\item
  对 A 的每个高度 1 的素理想 p，环 A \(_{p}\) 整闭；
\item
  对 A 的每个素理想 \(\mathfrak{p}\) （高度 \(\geqslant 2\) 者），有 \(\operatorname{prof}(A_{\mathfrak{p}}) \geqslant 2\) 。
\end{enumerate}

则 A 正规。

目标是证明环 \(A_{\mathfrak{p}}\) 对每个 \(\mathfrak{p} \in \operatorname{Spec}(A)\) 整闭，我们对 \(\operatorname{ht}(\mathfrak{p})\) 作归纳来证明这一点。对 \(\operatorname{ht}(\mathfrak{p}) \leqslant 1\)，这由假设 与 得出。于是设 \(\operatorname{ht}(\mathfrak{p}) \geqslant 2\) ，且 \(A_{\mathfrak{q}}\) 对 A 的每个素理想 \(\mathfrak{q}\) （满足 \(\operatorname{ht}(\mathfrak{q}) < \operatorname{ht}(\mathfrak{p})\) 者）整闭。由假设，有 \(\operatorname{prof}(A_{\mathfrak{p}}) \geqslant 2\) 。由归纳假设及第 9 节的定理 3 的推论 3——应用于环 \(A_{\mathfrak{p}}\) 及闭子集 \(\{\mathfrak{p}A_{\mathfrak{p}}\}\) （ \(\operatorname{Spec}(A_{\mathfrak{p}})\) 的闭子集）——环 \(A_{\mathfrak{p}}\) 是整环。设 K 为它的分式域，设 B 为 K 的一个包含 \(A_{\mathfrak{p}}\) 且在 \(A_{\mathfrak{p}}\) 上有限的子环。剩下要证明 B 等于 \(A_{\mathfrak{p}}\) 。用 i 表示从 \(A_{\mathfrak{p}}\) 到 B 的典范单射。由于 B 包含于 K，它是无挠的 \(A_{\mathfrak{p}}\) -模，故深度 \(\geqslant 1\) 。此外，对每个素理想 \(\mathfrak{q}\) （ \(A_{\mathfrak{p}}\) 的素理想，异于 \(\mathfrak{p}A_{\mathfrak{p}}\) 者），同态 \(i_{\mathfrak{q}}: (A_{\mathfrak{p}})_{\mathfrak{q}} \to B_{\mathfrak{q}}\) 是双射，因为 \(A_{\mathfrak{q}}\) 整闭。由第 9 节的引理 4——应用于闭子集 \(F = \{\mathfrak{p}A_{\mathfrak{p}}\}\) （ \(\operatorname{Spec}(A_{\mathfrak{p}})\) 的闭子集）——同态 i 是双射，这就完成了定理的证明。

注 1.--- 定理 4 的一个方便的形式如下：设 A 为这样的 Noether 环：对 A 的每个素理想 p，环 A \(_{p}\) 整闭或深度 ≥ 2；则 A 正规。事实上，我们来验证定理 4 的假设。设 p ∈ Spec(A)。若 ht(p) ≤ 1，则有 prof(A \(_{p}\) ) ≤ 1，故 A \(_{p}\) 整闭。若 ht(p) ≥ 2，则环 A \(_{p}\) 或者深度 ≥ 2，或者整闭，后者同样蕴含深度(A \(_{p}\) ) ≥ 2（第 1 节, 命题 1 的推论 2），由此得所要的结论。类似地可验证下述命题：设 A 为这样的 Noether 环：对 A 的每个素理想 p，环 A \(_{p}\) 约化或深度 ≥ 1；则 A 约化。

推论 1.--- 设 \(A\) 为 Noether 环，F 为 \(\operatorname{Spec}(A)\) 的闭子集，U 为其补开集。假设 \(\operatorname{prof}_{\mathrm{F}}(\mathrm{A})\) 为 \(\geqslant 2\) （相应地为 \(\geqslant 1\) ），且对每个 \(\mathfrak{p} \in \mathrm{U}\) ，环 \(\mathrm{A}_{\mathfrak{p}}\) 整闭（相应地约化）。则 A 正规（相应地约化）。

对每个 \(\mathfrak{p} \in \mathrm{F}\) ，有 \(\operatorname{prof}(\mathrm{A}_{\mathfrak{p}}) \geqslant \operatorname{prof}_{\mathrm{F}}(\mathrm{A})\) （ \(n^{\circ}\) 5, 命题 8）；因此只需应用前一个注。

推论 2.-- 设 ρ : A → B 为使 B 成为平坦 A-模的 Noether 环之间的同态。

\begin{enumerate}
\def\labelenumi{\alph{enumi})}
\item
  若 B 正规且在 A 上忠实平坦，则 A 正规。
\item
  设 A 正规，且环 \(\kappa(\mathfrak{p})\otimes_{\mathrm{A}}\mathrm{B}\) 当 \(\mathfrak{p}\) 为 A 的极小素理想时正规、当 \(\mathfrak{p}\) 为高度 1 的素理想时约化。则环 B 正规。
\item
  设 B 正规且在 A 上忠实平坦。则 \(\rho\) 是单射 (I, § 3, 第 5 节, 命题 9) 且 B 约化，故 A 约化。设 S 为 A 中非零因子构成的集合。由于 B 在 A 上平坦，\(\rho(S)\) 由 B 中的非零因子构成，故 B 在 \(\rho(S)^{-1}B\) 中整闭。设 \(x\) 为 \(S^{-1}A\) 中在 A 上整的元素；则元素 \(x \otimes 1_B\) ——环 \(S^{-1}A \otimes_A B\) （它等同于 \(\rho(S)^{-1}B\) ）的元素——在 B 上整，故属于 B。由于 B 在 A 上忠实平坦，有 \(x \in A\) (I, § 3, 第 5 节, 命题 10, (ii))，故 A 正规。
\item
  在 b) 的假设下，由注 1 只需证明：对 B 的每个素理想 \(\mathfrak{q}\) ，局部环 \(B_{\mathfrak{q}}\) 正规或深度 \(\geqslant 2\) 。用 \(\mathfrak{p}\) 表示 A 的素理想 \(\rho^{-1}(\mathfrak{q})\) ；局部同态 \(A_{\mathfrak{p}} \to B_{\mathfrak{q}}\) ——由 \(\rho\) 诱导——使 \(B_{\mathfrak{q}}\) 成为忠实平坦的 \(A_{\mathfrak{p}}\) -模 (I, § 3, 第 5 节, 命题 9)。我们区分四种情形：
\end{enumerate}

\begin{enumerate}
\def\labelenumi{\arabic{enumi})}
\item
  若 \(\operatorname{ht}(\mathfrak{p}) = 0\) ，则 \(A_{\mathfrak{p}}\) 是域，等于 \(\kappa(\mathfrak{p})\) ；环 \(A_{\mathfrak{p}} \otimes_{A} B\) 按假设正规。\(B_{\mathfrak{q}}\) 亦然，它是其一个分式环。
\item
  若 \(\operatorname{ht}(\mathfrak{p}) = 1\) 且 \(\operatorname{ht}(\mathfrak{q}) \leqslant 1\) ，则 \(A_{\mathfrak{p}}\) 是离散赋值环；设 \(\pi\) 为 \(A_{\mathfrak{p}}\) 的一个单值化元。由于 \(B_{\mathfrak{q}}\) 在 \(A_{\mathfrak{p}}\) 上忠实平坦，元素 \(\pi 1_{B_{\mathfrak{q}}}\) ——\(B_{\mathfrak{q}}\) 的元素——不是零因子，故局部环 \(B_{\mathfrak{q}}/\pi B_{\mathfrak{q}}\) 的维数为 0 (VIII, § 3, 第 1 节, 命题 1 的推论 2)。它是约化的，因为它是约化环 \(\kappa(\mathfrak{p}) \otimes_{A} B\) 的一个分式环，从而它是域。因此 \(B_{\mathfrak{q}}\) 是离散赋值环 (VI, § 1, 第 4 节, 命题 2)，故整闭。
\item
  \(\operatorname{ht}(\mathfrak{p}) = 1\) 且 \(\operatorname{ht}(\mathfrak{q}) \geqslant 2\) 。则由关系式
\end{enumerate}

\[
\dim (B _ {\mathfrak {q}}) = \dim (A _ {\mathfrak {p}}) + \dim (\kappa (\mathfrak {p}) \otimes_ {A} B _ {\mathfrak {q}})
\]

(VIII, § 3, 第 4 节, 命题 7 的推论 1)，环 \(\kappa(\mathfrak{p})\otimes_{\mathrm{A}}\mathrm{B}_{\mathfrak{q}}\) 的维数 \(\geqslant 1\) 。它按假设约化，故深度 \(\geqslant 1\) （第 1 节, 注 3）。于是有（第 6 节, 命题 11 的推论）

\[
\operatorname{prof} \left(\mathrm{B} _ {\mathfrak {q}}\right) = \operatorname{prof} \left(\Lambda_ {\mathfrak {p}}\right) + \operatorname{prof} \left(\kappa (\mathfrak {p}) \otimes_ {\mathrm{A}} \mathrm{B} _ {\mathfrak {q}}\right) \geqslant 1 + 1 = 2.
\]

\begin{enumerate}
\def\labelenumi{\arabic{enumi})}
\setcounter{enumi}{3}
\tightlist
\item
  \(\operatorname{ht}(\mathfrak{p}) \geqslant 2\) 。由于 \(A_{p}\) 的深度 \(\geqslant 2\) （命题 16 的推论），\(B_{q}\) 由同上引文亦然。
\end{enumerate}

推论 3.--- 设 \(\rho: A \to B\) 为 Noether 环之间的同态。设 \(B\) 是平坦 A-模，\(A\) 正规，且 \(\kappa(\mathfrak{p}) \otimes_A B\) 对所有 \(\mathfrak{p} \in \operatorname{Spec}(A)\) 正规。则 \(B\) 正规。

\section{Macaulay 模与环}

\subsection{Macaulay 模}

设 A 为 Noether 环，M 为有限生成 A-模，p 为 A 的素理想。若 \(p \notin \text{Supp}(M)\) ，则 \(M_{p} = 0\) ，从而 \(\text{prof}_{A_{p}}(M_{p}) = +\infty\) 且 \(\dim_{A_{p}}(M_{p}) = -\infty\) 。若 \(p \in \text{Supp}(M)\) ，则有 \(0 \leqslant \text{prof}_{A_{p}}(M_{p}) \leqslant \dim_{A_{p}}(M_{p}) < +\infty\) （§ 1, 第 4 节, 定理 2 的推论 2）。

定义 1.--- 设 A 为 Noether 环，M 为有限型的 \(\Lambda\) 模。称 M 是 Macaulay 的或为 Macaulay 模，如果对每个极大理想 \(\mathfrak{m} \in \operatorname{Supp}(M)\) ，有 \(\operatorname{prof}_{\mathrm{A}_{\mathfrak{m}}}\left(\mathrm{M}_{\mathfrak{m}}\right) = \dim_{\mathrm{A}_{\mathfrak{m}}}\left(\mathrm{M}_{\mathfrak{m}}\right)\) 。

由前所述，等价的说法是：有 \(\mathrm{prof}_{\mathrm{A}_{\mathfrak{m}}}\left(\mathrm{M}_{\mathfrak{m}}\right) \geqslant \dim_{\mathrm{A}_{\mathfrak{m}}}\left(\mathrm{M}_{\mathfrak{m}}\right)\) 对 A 的每个极大理想 \(\mathfrak{m}\) 成立。设 A 为 Noether 局部环；一个非零的有限生成 A-模是 Macaulay 的，必须且只需其深度等于其维数。

例.--- 1) 每个有限长度的 A-模都是 Macaulay 的。

\begin{enumerate}
\def\labelenumi{\arabic{enumi})}
\setcounter{enumi}{1}
\tightlist
\item
  设 M' 为有限生成 Macaulay A-模 M 的一个直和项子模。则 M' 是 Macaulay 的；事实上，对 A 的每个极大理想 m，\(A_{m}\) -模 \(M'_{m}\) 是 \(M_{m}\) 的直和项，因此有
\end{enumerate}

\[
\operatorname{prof} _ {A _ {\mathfrak {m}}} \left(M _ {\mathfrak {m}} ^ {\prime}\right) \geqslant \operatorname{prof} _ {A _ {\mathfrak {m}}} \left(M _ {\mathfrak {m}}\right) \geqslant \dim_ {A _ {\mathfrak {m}}} \left(M _ {\mathfrak {m}}\right) \geqslant \dim_ {A _ {\mathfrak {m}}} \left(M _ {\mathfrak {m}} ^ {\prime}\right),
\]

由 § 1, 第 1 节的注 4 及 VIII, § 1, 第 4 节, 命题 9 c)。

\begin{enumerate}
\def\labelenumi{\arabic{enumi})}
\setcounter{enumi}{2}
\tightlist
\item
  设 M 为 A 上的有限生成 Macaulay 模，\((x_{1},\ldots,x_{n})\) 为由 \(\Lambda\) 中元素组成的 M-正则序列。则 A-模 \(\overline{\mathrm{M}}=\mathrm{M}/(x_{1}\mathrm{M}+\cdots+x_{n}\mathrm{M})\) 是 Macaulay 的。事实上，设 m 为属于 \(\overline{M}\) 的支撑集的 A 的一个极大理想；有 \(x_{i}\in m\) 对每个 i 成立，因为 \(x_{i}\) 零化 \(\overline{M}\) ，且 \(x_{i}\) 在 \(A_{m}\) 中的典范像组成一个 \(M_{m}\) -正则序列，由 \(mA_{m}\) 中元素组成。因此（§ 1, 第 4 节, 命题 7 及 VIII, § 3, 第 2 节, 命题 3 的推论）有等式
\end{enumerate}

\[
\begin{array}{l} \operatorname{prof} _ {\mathrm{A} _ {\mathfrak {m}}} (\overline {{\mathrm{M}}} _ {\mathfrak {m}}) = \operatorname{prof} _ {\mathrm{A} _ {\mathfrak {m}}} (\mathrm{M} _ {\mathfrak {m}}) - n, \\ \dim_ {\mathrm{A} _ {\mathfrak {m}}} (\overline {{\mathrm{M}}} _ {\mathfrak {m}}) = \dim_ {\mathrm{A} _ {\mathfrak {m}}} (\mathrm{M} _ {\mathfrak {m}}) - n, \end{array}
\]

由此得我们的断言。

\begin{enumerate}
\def\labelenumi{\arabic{enumi})}
\setcounter{enumi}{3}
\tightlist
\item
  设 M 为有限生成 A-模，a 为 A 的理想且 \(aM = 0\) 。A-模 M 是 Macaulay 的，必须且只需它作为 \((A/a)\) -模是 Macaulay 的。事实上，令 B = A/a；设 n 为 B 的极大理想，m 为它在 \(\Lambda\) 中的原像；有 \(\dim_{\mathrm{A}_{m}}(\mathrm{M}_{\mathfrak{m}}) = \dim_{\mathrm{B}_{n}}(\mathrm{M}_{\mathfrak{n}})\) 及 \(\operatorname{prof}_{\mathrm{A}_{m}}(\mathrm{M}_{\mathfrak{m}}) = \operatorname{prof}_{\mathrm{B}_{n}}(\mathrm{M}_{\mathfrak{n}})\) （§ 1, 第 3 节, 命题 4 的推论）。
\end{enumerate}

命题 1.--- 设 \(\Lambda\) 为 Noether 环，M 为有限生成 A-模，\(\mathfrak{p}\) 与 \(\mathfrak{q}\) 为 Supp(M) 的素理想且 \(\mathfrak{p} \subset \mathfrak{q}\) 。设 \(\dim_{\mathrm{A}_{\mathfrak{q}}}(\mathrm{M}_{\mathfrak{q}}) = \operatorname{prof}_{\mathrm{A}_{\mathfrak{q}}}(\mathrm{M}_{\mathfrak{q}})\) 。则有 \(\dim_{\mathrm{A}_{\mathfrak{p}}}(\mathrm{M}_{\mathfrak{p}}) = \operatorname{prof}_{\mathrm{A}_{\mathfrak{p}}}(\mathrm{M}_{\mathfrak{p}})\) 以及

\[
\dim_ {A _ {q}} (M _ {q}) = \dim_ {A _ {p}} (M _ {p}) + \dim (A _ {q} / p A _ {q}).
\]

这直接由 § 1, 第 7 节的命题 13 的推论 1 得出。

推论.--- 设 A 为 Noether 环，M 为有限生成 A-模。下列条件等价：

\begin{enumerate}
\def\labelenumi{(\roman{enumi})}
\item
  A-模 M 是 Macaulay 的
\item
  有 \(a \operatorname{prof}_{\mathrm{A}_{\mathfrak{p}}}(\mathrm{M}_{\mathfrak{p}}) = \dim_{\mathrm{A}_{\mathfrak{p}}}(\mathrm{M}_{\mathfrak{p}})\) 对每个 \(\mathfrak{p} \in \operatorname{Supp}(\mathrm{M})\) 成立；
\item
  有 \(\operatorname{prof}_{\mathrm{F}}(\mathrm{M}) = \operatorname{codim}(\operatorname{Supp}(\mathrm{M}) \cap \mathrm{F}, \operatorname{Supp}(\mathrm{M}))\) ，其中 \(\mathrm{F}\) 取遍 \(\operatorname{Spec}(\mathrm{A})\) 的闭子集；
\item
  有 \(a \operatorname{prof}_{\mathrm{A}}(\mathfrak{p}; \mathrm{M}) = \dim_{\mathrm{A}_{\mathfrak{p}}}(\mathrm{M}_{\mathfrak{p}})\) 对每个 \(\mathfrak{p} \in \operatorname{Supp}(\mathrm{M})\) 成立。
\item
  \(\Rightarrow\) (ii)：由命题 1 得出。
\item
  \(\Rightarrow\) (iii)：由 § 1, 第 5 节的命题 8，\(\operatorname{prof}_{\mathrm{F}}(\mathbf{M})\) 是整数 \(\operatorname{prof}(\mathbf{M}_{\mathfrak{p}})\) 当 \(\mathfrak{p}\) 取遍 \(\operatorname{Supp}(\mathbf{M}) \cap \mathbf{F}\) 时的下确界。若 M 是 Macaulay 的，则对这样的 \(\mathfrak{p}\) 有等式 \(\operatorname{prof}(\mathbf{M}_{\mathfrak{p}}) = \dim(\mathbf{M}_{\mathfrak{p}}) = \operatorname{codim}(\mathbf{V}(\mathfrak{p}), \operatorname{Supp}(\mathbf{M}))\) (VIII, § 1, 第 4 节, 命题 9)，由此得。
\item
  \(\Rightarrow\) (iv)：只需取 \(F = V(\mathfrak{p})\) 。
\item
  \(\Rightarrow\) (i)：这由不等式 \(\operatorname{prof}_{\mathbf{A}}(\mathfrak{p};\mathbf{M})\leqslant\operatorname{prof}(\mathbf{M}_{\mathfrak{p}})\leqslant\dim(\mathbf{M}_{\mathfrak{p}})\) 得出，它对每个 \(\mathfrak{p}\in\operatorname{Supp}(\mathbf{M})\) 成立（ \(\S 1,\ n^{\circ}5\) , 注 3 及 \(n^{\circ}4\) , 定理 2 的推论 2）。
\end{enumerate}

注.--- 设 S 为 A 的乘性子集，M 为有限生成 Macaulay A-模。则 \(S^{-1}M\) 是 Macaulay 的 \(S^{-1}A\) -模。事实上，设 \(\mathfrak{q} \in \operatorname{Spec}(S^{-1}A)\)；用 \(i_{\mathrm{A}}^{\mathrm{S}} : A \to S^{-1}A\) 表示典范同态，并令 \(\mathfrak{p} = (i_{\mathrm{A}}^{\mathrm{S}})^{-1}(\mathfrak{q})\)。环 \((S^{-1}A)_{\mathfrak{q}}\) 等同于 \(A_{\mathfrak{p}}\) (II, § 2, 第 5 节, 命题 11)，且 \(A_{\mathfrak{p}}\) -模 \((S^{-1}M)_{\mathfrak{q}}\) 等同于 \(A_{\mathfrak{p}}\) -模 \(M_{\mathfrak{p}}\) (II, § 2, 第 7 节, 命题 20)，后者由推论是 Macaulay 的。

\subsection{Macaulay 模的支撑集}

命题 2. — 设 A 为 Noether 环，M 为有限生成 Macaulay A-模。

\begin{enumerate}
\def\labelenumi{\alph{enumi})}
\item
  A-模 M 没有嵌入相伴素理想。 \(^{1}\)
\item
  设 X 为 Supp(M) 的闭不可约子集，Y 为 X 的闭子集。有
\end{enumerate}

\[
\operatorname{codim} (Y, X) + \operatorname{codim} (X, \operatorname{Supp} (M)) = \operatorname{codim} (Y, \operatorname{Supp} (M)).
\]

\begin{enumerate}
\def\labelenumi{\alph{enumi})}
\setcounter{enumi}{2}
\item
  拓扑空间 Supp(M) 是链式的 (VIII, § 1, 第 2 节, 定义 4)。
\item
  设 \(X_{1}\) 与 \(X_{2}\) 为 Supp(M) 的不可约分支，Y 为 \(X_{1} \cap X_{2}\) 的闭子集。有 codim(Y, \(X_{1}\) ) = codim(Y, \(X_{2}\) )。
\item
  设 \(\mathfrak{p} \in \operatorname{Ass}(M)\) 。有 \(\operatorname{prof}_{\mathrm{A}}(\mathfrak{p}; M) = 0\) （ \(\S\) 1, \(n^{\circ}\) 1, 注 2），从而 \(\dim(M_{\mathfrak{p}}) = 0\) （ \(n^{\circ}\) 1, 命题 1 的推论），这蕴含 \(\mathfrak{p}\) 是 \(\operatorname{Supp}(M)\) 的极小元。
\item
  先设 Y 不可约。设 p 与 q 为 Supp(M) 的素理想，使得 \(Y = V(\mathfrak{q})\) 且 \(X = V(\mathfrak{p})\) 。由命题 1 我们有
\end{enumerate}

\[
\begin{array}{r l} \operatorname{codim} (Y, X) & = \dim (A _ {q} / p A _ {q}) = \dim (M _ {q}) - \dim (M _ {p}) \\ & = \operatorname{codim} (Y, \operatorname{Supp} (M)) - \operatorname{codim} (X, \operatorname{Supp} (M)). \end{array}
\]

一般情形由 VIII, § 1, 第 2 节, 注 3 得出。

\begin{enumerate}
\def\labelenumi{\alph{enumi})}
\setcounter{enumi}{2}
\tightlist
\item
  设 X、Y、Z 为 Supp(M) 的闭不可约子集且 \(Z \subset Y \subset X\) 。这些子集在 Supp(M) 中的余维数均有限 (VIII, § 1, 第 4 节, 命题 9 及 § 3, 第 1 节, 命题 2 的推论 1)。于是由 b) 推出等式
\end{enumerate}

\[
\operatorname{codim} (Z, Y) + \operatorname{codim} (Y, X) = \operatorname{codim} (Z, X)
\]

它蕴含 c) (VIII, § 1, 第 2 节, 命题 4)。

\begin{enumerate}
\def\labelenumi{\alph{enumi})}
\setcounter{enumi}{3}
\tightlist
\item
  由 b)，有 \(\text{codim}(Y, X_1) = \text{codim}(Y, \text{Supp}(M)) = \text{codim}(Y, X_2)\) 。
\end{enumerate}

特别地，若存在支撑集等于 Spec(A) 的有限生成 Macaulay A-模 M，则环 A 是链式的，从而 A 的每个分式环或每个商环都是链式的 (VIII, § 1, 第 3 节, 注 2)。

注 1. — 在命题 2 的假设下，可能发生这样的情形：Supp(M) 的两个不可约分支 \(X_{1}\) 与 \(X_{2}\) 的交 Y 退化为一点，且有 \(\dim X_{1} \neq \dim X_{2}\) 及 \(\dim X_{2} \neq \text{codim}(Y, X_{2})\) （参见习题 4）。然而，当环 A 是局部环时这不可能发生，如下面的推论所示。

推论. — 设 A 为 Noether 局部环，M 为非零的有限生成 Macaulay A-模。

\begin{enumerate}
\def\labelenumi{\alph{enumi})}
\item
  Supp(M) 的闭不可约子集的所有极大链长度都等于 dim(M)。
\item
  对 Supp(M) 的每个闭子集 X，有
\end{enumerate}

\[
\operatorname{codim} (\mathrm{X}, \operatorname{Supp} (\mathrm{M})) = \dim (\operatorname{Supp} (\mathrm{M})) - \dim (\mathrm{X}).
\]

\begin{enumerate}
\def\labelenumi{\alph{enumi})}
\setcounter{enumi}{2}
\item
  Supp(M) 的所有不可约分支维数相同。
\item
  对 A 的每个理想 J，有
\end{enumerate}

\[
\operatorname{prof} _ {\mathrm{A}} (\mathrm{J}; \mathrm{M}) = \dim (\mathrm{M}) - \dim (\mathrm{M} / \mathrm{JM}).
\]

\begin{enumerate}
\def\labelenumi{\alph{enumi})}
\tightlist
\item
  Supp(M) 的闭不可约子集的一个极大链，其最小元为 \(\{\mathfrak{m}_{\mathrm{A}}\}\) ，其最大元为不可约分支 X，X 属于
\end{enumerate}

Supp(M)。它的长度等于 \(\{m_{A}\}\) 在 X 中的余维数（命题 2, c)）；由命题 2, b) 应用于闭子集 \(\{m_{A}\}\subset X\) ，它等于 \(\text{codim}(\{\mathfrak{m}_{\Lambda}\},\text{Supp}(M))\) ，即 \(\dim(M)\) 。

\begin{enumerate}
\def\labelenumi{\alph{enumi})}
\setcounter{enumi}{1}
\item
  当子集 X 不可约时这是 a) 的推论 (VIII, § 1, 第 2 节, 命题 5)；一般情形由 VIII, § 1, 第 1 节, 命题 1 及 § 1, 第 2 节, 注 4 得出。
\item
  这是 b) 的推论。
\item
  有 \(\operatorname{prof}_{\mathrm{A}}(\mathrm{J};\mathrm{M}) = \operatorname{codim}(\operatorname{Supp}(\mathrm{M})\cap \mathrm{V}(\mathrm{J}),\operatorname{Supp}(\mathrm{M}))\) ，由第 1 节的命题 1 的推论。于是只需取 \(\mathrm{X} = \operatorname{Supp}(\mathrm{M})\cap \mathrm{V}(\mathrm{J}) = \operatorname{Supp}(\mathrm{M} / \mathrm{JM})\) 应用 b) (II, § 4, 第 4 节, 命题 18 的推论)。
\end{enumerate}

注 2.--- 设 M 为有限生成 Macaulay A-模，p 为 Supp(M) 的元素。由 § 1, 第 4 节的定理 2，有 prof\_A(p; M) \textless{} +∞，且存在长度为 prof\_A(p; M) 的、由 p 中元素组成的 M-正则序列。用 J 表示由这样一个序列生成的 A 的理想；则 A-模 M/JM 是 Macaulay 的（第 1 节, 例 3），且 p 是其支撑集的极小元。事实上，p 包含 J，从而属于 M/JM 的支撑集 (II, § 4, 第 4 节, 命题 18 的推论)；由 § 1, 第 4 节的定理 2 的推论 1，理想 p 包含于 Ass(M/JM) 的某个元素，但与有限生成 Macaulay 模相伴的每个素理想都是其支撑集的极小元（命题 2）。
